% Options for packages loaded elsewhere
\PassOptionsToPackage{unicode}{hyperref}
\PassOptionsToPackage{hyphens}{url}
\PassOptionsToPackage{dvipsnames,svgnames,x11names}{xcolor}
%
\documentclass[
  10pt,
]{article}
\usepackage{amsmath,amssymb}
\usepackage{iftex}
\ifPDFTeX
  \usepackage[T1]{fontenc}
  \usepackage[utf8]{inputenc}
  \usepackage{textcomp} % provide euro and other symbols
\else % if luatex or xetex
  \usepackage{unicode-math} % this also loads fontspec
  \defaultfontfeatures{Scale=MatchLowercase}
  \defaultfontfeatures[\rmfamily]{Ligatures=TeX,Scale=1}
\fi
\usepackage{lmodern}
\ifPDFTeX\else
  % xetex/luatex font selection
  \setmainfont[]{Latin Modern Roman}
  \setmonofont[]{DejaVu Sans Mono}
\fi
% Use upquote if available, for straight quotes in verbatim environments
\IfFileExists{upquote.sty}{\usepackage{upquote}}{}
\IfFileExists{microtype.sty}{% use microtype if available
  \usepackage[]{microtype}
  \UseMicrotypeSet[protrusion]{basicmath} % disable protrusion for tt fonts
}{}
\makeatletter
\@ifundefined{KOMAClassName}{% if non-KOMA class
  \IfFileExists{parskip.sty}{%
    \usepackage{parskip}
  }{% else
    \setlength{\parindent}{0pt}
    \setlength{\parskip}{6pt plus 2pt minus 1pt}}
}{% if KOMA class
  \KOMAoptions{parskip=half}}
\makeatother
\usepackage{xcolor}
\usepackage[margin=2.3cm]{geometry}
\usepackage{longtable,booktabs,array}
\usepackage{calc} % for calculating minipage widths
% Correct order of tables after \paragraph or \subparagraph
\usepackage{etoolbox}
\makeatletter
\patchcmd\longtable{\par}{\if@noskipsec\mbox{}\fi\par}{}{}
\makeatother
% Allow footnotes in longtable head/foot
\IfFileExists{footnotehyper.sty}{\usepackage{footnotehyper}}{\usepackage{footnote}}
\makesavenoteenv{longtable}
\setlength{\emergencystretch}{3em} % prevent overfull lines
\providecommand{\tightlist}{%
  \setlength{\itemsep}{0pt}\setlength{\parskip}{0pt}}
\setcounter{secnumdepth}{-\maxdimen} % remove section numbering
\usepackage{mathtools}
\renewcommand{\contentsname}{Sumário}
\ifLuaTeX
  \usepackage{selnolig}  % disable illegal ligatures
\fi
\IfFileExists{bookmark.sty}{\usepackage{bookmark}}{\usepackage{hyperref}}
\IfFileExists{xurl.sty}{\usepackage{xurl}}{} % add URL line breaks if available
\urlstyle{same}
\hypersetup{
  colorlinks=true,
  linkcolor={blue},
  filecolor={Maroon},
  citecolor={Blue},
  urlcolor={Blue},
  pdfcreator={LaTeX via pandoc}}

\author{}
\date{}

\begin{document}

\noindent\textbf{Controle normativo P10-A.}
As expressões operacionais de capacidade e demanda são regidas por
\texttt{CAPACIDADE\_E\_RECURSOS.md}: \(q_{ij}\le
F_i(m_{ij})h_id_i\eta_{p,j}\eta_{w,j}\), sem fator de ineficiência
\(\eta_s\). As expressões financeiras históricas baseadas em
\(\gamma/\rho\) permanecem condicionadas à revisão da issue \#11
e não podem, isoladamente, sustentar certificação GO.


% P10-A NORMATIVE: docs/product/5-specifications/formulario_longo/CAPACIDADE_E_RECURSOS.md
% q_ij = F_i(m_ij) * h_i * d_i * eta_p,j * eta_w,j (capped by remaining Q).
% eta_s is NOT an allocation inefficiency. Old rho*eta_s*m and derived
% costs are non-certifiable for nonlinear F pending issue #11.


% P09-A: common normative contract for the MD, summary and TeX representations.
% See docs/product/5-specifications/formulario_longo/CONVENCOES_NORMATIVAS.md.
% Units, ramp domains, half-open contract intervals and economic tie-breaks
% are normative there; equations under #10--#13 remain subject to proof.


{
\hypersetup{linkcolor=}
\setcounter{tocdepth}{4}
\tableofcontents
}
\hypertarget{formuluxe1rio-white-box-atividades-pool-de-recursos-e-custos-indiretos}{%
\section{Formulário White-Box --- Atividades, Pool de Recursos e Custos
Indiretos}\label{formuluxe1rio-white-box-atividades-pool-de-recursos-e-custos-indiretos}}

\begin{center}\rule{0.5\linewidth}{0.5pt}\end{center}

\hypertarget{parte-i-entrada-definiuxe7uxe3o-das-atividades-formuluxe1rio}{%
\subsection{Parte I --- Entrada: Definição das Atividades
(formulário)}\label{parte-i-entrada-definiuxe7uxe3o-das-atividades-formuluxe1rio}}

\hypertarget{i.1-conjunto-de-atividades}{%
\subsubsection{I.1 Conjunto de
atividades}\label{i.1-conjunto-de-atividades}}

\[\mathcal{A} = \{A_i\}_{i=1}^{n_{\text{ativ}}}\]

\hypertarget{i.2-tupla-da-atividade}{%
\subsubsection{I.2 Tupla da atividade}\label{i.2-tupla-da-atividade}}

Na prática, cada atividade representa uma frente de trabalho do projeto
--- escavação, concretagem, lançamento de base, e assim por diante ---
com sua própria meta física, sua equipe e suas dependências de outras
atividades. Os quatro campos abaixo resumem tudo que precisa ser
fornecido; os demais campos da tabela são calculados automaticamente a
partir deles.

\[A_i = (Q_i,\ E_i,\ W_i,\ B_i)\]

\begin{longtable}[]{@{}
  >{\raggedright\arraybackslash}p{(\columnwidth - 4\tabcolsep) * \real{0.3333}}
  >{\raggedright\arraybackslash}p{(\columnwidth - 4\tabcolsep) * \real{0.3333}}
  >{\raggedright\arraybackslash}p{(\columnwidth - 4\tabcolsep) * \real{0.3333}}@{}}
\toprule\noalign{}
\begin{minipage}[b]{\linewidth}\raggedright
Campo
\end{minipage} & \begin{minipage}[b]{\linewidth}\raggedright
Tipo
\end{minipage} & \begin{minipage}[b]{\linewidth}\raggedright
Preenchido por
\end{minipage} \\
\midrule\noalign{}
\endhead
\bottomrule\noalign{}
\endlastfoot
\(Q_i\) & quantitativo total {[}un{]} & usuário \\
\(E_i\) & equipe (§I.3) & usuário \\
\(W_i\) & política de mobilização (§I.4) & usuário \\
\(B_i\) & predecessores + completude parcial (§I.5) & usuário \\
\(j_i^0\) & semana de início & \textbf{derivado} (schedule, §I.5) \\
\(N_i\) & duração em semanas & \textbf{derivado} (fecho de escopo exato
via busca monotônica, Parte IV/XIV) \\
\((h_i, d_i)\) & regime de jornada da equipe & usuário,
\(d_i\in\{4,5,6\}\), \(h_i\in\{6,8\}\) \\
\end{longtable}

\hypertarget{i.3-equipe-e_i-mathcalg_i-r_i}{%
\subsubsection{\texorpdfstring{I.3 Equipe
\(E_i = (\mathcal{G}_i, R_i)\)}{I.3 Equipe E\_i = (\textbackslash mathcal\{G\}\_i, R\_i)}}\label{i.3-equipe-e_i-mathcalg_i-r_i}}

\hypertarget{i.3.1-composiuxe7uxe3o-de-recursos}{%
\paragraph{I.3.1 Composição de
recursos}\label{i.3.1-composiuxe7uxe3o-de-recursos}}

\(R_i = \{r\}\), cada recurso descrito por:

\begin{longtable}[]{@{}
  >{\raggedright\arraybackslash}p{(\columnwidth - 2\tabcolsep) * \real{0.5000}}
  >{\raggedright\arraybackslash}p{(\columnwidth - 2\tabcolsep) * \real{0.5000}}@{}}
\toprule\noalign{}
\begin{minipage}[b]{\linewidth}\raggedright
Campo
\end{minipage} & \begin{minipage}[b]{\linewidth}\raggedright
Significado
\end{minipage} \\
\midrule\noalign{}
\endhead
\bottomrule\noalign{}
\endlastfoot
\(\text{cat}(r) \in \{EQ, MO, MAT, COMB\}\) & categoria de BDI \\
\(n_r\) & quantidade do recurso \(r\) dentro de \textbf{uma} frente \\
regime & \texttt{mensal} (EQ/MO) ou \texttt{consumo} (MAT/COMB) \\
\(V_r^{\text{mensal}}\) & custo contratual mensal {[}R\$/mês{]}, se
regime = \texttt{mensal} \\
\(\kappa_r\) & taxa física de consumo
{[}unid\(_g\)/unid\(_{\text{tempo}}\){]}, se regime =
\texttt{consumo} \\
\(g(r)\) & tipo físico do bem consumível, se regime =
\texttt{consumo} \\
\(p_{g(r)}\) & preço de aquisição {[}R\$/unid\(_g\){]}, se regime =
\texttt{consumo} \\
\end{longtable}

\textbf{Custo horário de catálogo do recurso \(r\)} --- agnóstico em
unidade física:

\[c_r = \begin{cases}
\dfrac{V_r^{\text{mensal}}}{s_{\text{mes}}\cdot d_i\cdot h_i} & \text{cat}(r)\in\{EQ,MO\}\quad\text{(diluição contratual)}\\[10pt]
\kappa_r \cdot p_{g(r)} & \text{cat}(r)\in\{MAT,COMB\}\quad\text{(consumo: taxa física × preço unitário)}
\end{cases}\]

onde: - \(\kappa_r\) {[}unid\(_g\) / h\(_{\text{op}}\){]}: taxa de
consumo do bem \(g(r)\) por hora de operação do recurso \(r\) -
\(p_{g(r)}\) {[}R\$/unid\(_g\){]}: preço de aquisição por unidade física
do bem \(g\) - \textbf{unid\(_g\)} é declarada no cadastro do bem e não
é fixada neste formulário

\emph{Exemplos de instâncias da mesma fórmula:} combustível →
\(\kappa_r\) {[}L/h{]}, \(p_g\) {[}R\$/L{]}; agregado → \(\kappa_r\)
{[}t/h{]}, \(p_g\) {[}R\$/t{]}; explosivo → \(\kappa_r\) {[}kg/h{]},
\(p_g\) {[}R\$/kg{]}; cimento → \(\kappa_r\) {[}kg/h{]}, \(p_g\)
{[}R\$/kg{]}.

\textbf{MAT legado.} Quando
\texttt{atributos\_especificos.consumo\_horario\_fisico} e
\texttt{atributos\_especificos.unidade\_fisica} estão ausentes no
cadastro do bem, trata-se do caso degenerado
\(\kappa_r = 1\ [\text{unid}_g/\text{h}]\),
\(p_{g(r)} = c_r\ [\text{R\$}/\text{h}]\) --- a fórmula
\(c_r = \kappa_r \cdot p_{g(r)}\) permanece válida com a unidade física
colapsada. Quando os atributos estão presentes, a decomposição é
auditável por tipo físico, retrocompatível com catálogos que não
distinguem taxa de consumo e preço unitário.

\hypertarget{i.3.2-grafo-de-sinergia-e-produtividade-32-da-referuxeancia}{%
\paragraph{I.3.2 Grafo de sinergia e produtividade (§32 da
referência)}\label{i.3.2-grafo-de-sinergia-e-produtividade-32-da-referuxeancia}}

O \textbf{grafo de sinergia}
\(\mathcal{G}_i = (V_i,\ A_i^{\text{syn}},\ \text{cap})\) é um dígrafo
ponderado sobre os recursos da equipe:

\begin{longtable}[]{@{}
  >{\raggedright\arraybackslash}p{(\columnwidth - 2\tabcolsep) * \real{0.5000}}
  >{\raggedright\arraybackslash}p{(\columnwidth - 2\tabcolsep) * \real{0.5000}}@{}}
\toprule\noalign{}
\begin{minipage}[b]{\linewidth}\raggedright
Elemento
\end{minipage} & \begin{minipage}[b]{\linewidth}\raggedright
Descrição
\end{minipage} \\
\midrule\noalign{}
\endhead
\bottomrule\noalign{}
\endlastfoot
\(V_i \supseteq R_i\) & nós: recursos + etapas intermediárias de
produção \\
\(A_i^{\text{syn}} \subseteq V_i \times V_i\) & arestas: dependências de
produção (fluxo de insumo para produto) \\
\(\text{cap}: A_i^{\text{syn}} \to \mathbb{R}_{>0}\) & capacidade da
aresta {[}unid\(_Q\) / unid\(_{\text{tempo}}\){]} \\
\(D \subseteq V_i\) & \textbf{domínio de produção}: subgrafo que
efetivamente gera a saída \(Q_i\) \\
\(s_i\) & fonte virtual (insumos brutos) \\
\(t_i\) & sorvedouro virtual (produto final na unidade de \(Q_i\)) \\
\end{longtable}

\textbf{Produtividade via min-corte sobre \(D\)} (teorema max-fluxo /
min-corte):

\[\boxed{\rho_i = \max\text{-fluxo}(s_i \to t_i,\ \mathcal{G}_i^D) = \min_{C \,\in\, \text{cortes}(s_i,t_i)} \sum_{(u,v)\in C} \text{cap}(u,v) \qquad \text{[unid}_Q\text{/h]}}\]

\textbf{Casos práticos} (topologias típicas de equipes de construção):

\begin{longtable}[]{@{}
  >{\raggedright\arraybackslash}p{(\columnwidth - 4\tabcolsep) * \real{0.3333}}
  >{\raggedright\arraybackslash}p{(\columnwidth - 4\tabcolsep) * \real{0.3333}}
  >{\raggedright\arraybackslash}p{(\columnwidth - 4\tabcolsep) * \real{0.3333}}@{}}
\toprule\noalign{}
\begin{minipage}[b]{\linewidth}\raggedright
Topologia
\end{minipage} & \begin{minipage}[b]{\linewidth}\raggedright
Estrutura de \(\mathcal{G}_i^D\)
\end{minipage} & \begin{minipage}[b]{\linewidth}\raggedright
Fórmula fechada
\end{minipage} \\
\midrule\noalign{}
\endhead
\bottomrule\noalign{}
\endlastfoot
\textbf{Série} (gargalo único) & \(r_1 \to r_2 \to \cdots \to r_k\) &
\(\rho_i = \min_{(u,v)\,\in\,\text{caminho}} \text{cap}(u,v)\) \\
\textbf{Paralelo} (aditivo) &
\(s_i \rightrightarrows \{r_1,\ldots,r_k\} \rightrightarrows t_i\) &
\(\rho_i = \sum_{r} \text{cap}(s_i,r)\) \\
\textbf{Misto} & combinação de série e paralelo & algoritmo
Ford-Fulkerson sobre \(\mathcal{G}_i^D\) \\
\end{longtable}

\textbf{Derivação de \(\rho_i\).} Quando o grafo de sinergia
\(\mathcal{G}_i\) está presente no cadastro da equipe, \(\rho_i\) é
derivado por \texttt{producao\_from\_grafo(G\_i)} via min-corte --- o
campo \texttt{producao\_horaria} no artefato \texttt{Producao} funciona
como cache auditável do resultado. Quando ausente, \(\rho_i\) é entrada
direta, calibrada externamente (ficha técnica). A calibração das
capacidades \(\text{cap}(\cdot)\) das arestas permanece sempre externa a
este formulário; a incerteza do modelo reside nas capacidades unitárias
declaradas, não na fórmula de derivação de \(\rho_i\) a partir delas.

\textbf{Mobilização e capacidade física.} O número de frentes
\(m_{ij}\) determina os recursos realmente necessários, não uma
eficiência artificial. A função de produção inclui todas as frentes e
gargalos compartilhados:

\[\boxed{F_i(m)=\max\text{-fluxo}(s_i,t_i;\mathcal G_i^D\mid_m)}\]
\[\boxed{B_{ij}=F_i(m_{ij})h_id_i\eta_{p,j}\eta_{w,j}}\]

A taxa \(\rho_i^{peak}=F_i(m_i^{max})\) não é produtividade de uma
frente. O eventual quociente \(F_i(m)/F_i(m_i^{max})\) é
\textbf{diagnóstico de capacidade}, não uma perda que se possa multiplicar
novamente. Demanda física \(d_{gj}=\sum_i a_{ig}(m_{ij})\), presença
\(x_{gj}\ge d_{gj}\) e standby \(s_{gj}=x_{gj}-d_{gj}\ge0\) são estados
separados. A formulação e as contraprovas estão no contrato P10-A.

\hypertarget{i.4-poluxedtica-de-mobilizauxe7uxe3o-w_i}{%
\subsubsection{\texorpdfstring{I.4 Política de mobilização
\(W_i\)}{I.4 Política de mobilização W\_i}}\label{i.4-poluxedtica-de-mobilizauxe7uxe3o-w_i}}

\[W_i = (m_i^0,\ \mu_i^+,\ m_i^1,\ \mu_i^-,\ m_i^{\max}) \in \mathbb{N}_{>0}\times\mathbb{N}_{\ge0}\times\mathbb{N}_{>0}\times\mathbb{N}_{\ge0}\times\mathbb{N}_{>0}\]

\begin{longtable}[]{@{}ll@{}}
\toprule\noalign{}
Símbolo & Significado \\
\midrule\noalign{}
\endhead
\bottomrule\noalign{}
\endlastfoot
\(m_i^0\) & nº de frentes na semana de início \\
\(\mu_i^+\) & taxa de ramp-up {[}frentes/semana{]} \\
\(m_i^{\max}\) & nº máximo de frentes concorrentes (platô) \\
\(\mu_i^-\) & taxa de ramp-down {[}frentes/semana{]} \\
\(m_i^1\) & nº de frentes na última semana (piso final, nunca zero) \\
\end{longtable}

\textbf{Perfil trapezoidal} --- mínimo de três retas, semana relativa
\(\tau = j - j_i^0 \in \{0,\ldots,N_i-1\}\):

\[\boxed{m_i(\tau) = \min\!\Big(m_i^0 + \mu_i^+\tau,\quad m_i^{\max},\quad m_i^1 + \mu_i^-(N_i-1-\tau)\Big)}\]

\textbf{Caso degenerado retangular} --- quando
\(\mu_i^+ = \mu_i^- = 0\):

\[\mu_i^+ = \mu_i^- = 0 \implies m_i(\tau) = m_i^{\max} \quad \forall\,\tau \in \{0,\ldots,N_i-1\}\]

com \(m_i^0 = m_i^1 = m_i^{\max}\), reduzindo \(W_i\) a um único
parâmetro de frota: frota constante durante toda a atividade.

\[m_{ij} = \begin{cases} m_i(j-j_i^0) & j\in[j_i^0,\,j_i^0+N_i-1] \\ 0 & \text{caso contrário}\end{cases}\]

\textbf{Determinação de \(N_i\).} Para perfis não-retangulares, \(N_i\)
não possui fórmula fechada. A busca que a resolve não é circular: para
cada duração candidata \(N\), o perfil \(m_i(\tau;N)\) acima é uma
função de \(\tau\) inteiramente determinada --- nada nela depende de um
\(N_i\) ainda desconhecido ---, de modo que a produção acumulada
\(\Pi(N)\equiv\sum_{\tau=0}^{N-1}F_i(m_i(\tau;N))\varepsilon_{i\tau}^{exog}(N)h_id_i\)
pode ser calculada exatamente para qualquer \(N\) testado. O que não se
conhece de antemão é apenas em qual \(N\) essa soma cruza \(Q_i\) --- a
mesma situação de qualquer grandeza definida como ``o menor \(N\) tal
que uma condição monotônica se satisfaz'', sem mistério adicional.

A validade da busca depende de \(\Pi\) ser monotônica em \(N\) --- e
essa monotonicidade precisa ser justificada com cuidado, porque \(N\)
aparece dentro do próprio perfil (no termo de ramp-down), de modo que
aumentar \(N\) não apenas acrescenta uma semana: redesenha o perfil
inteiro. Ainda assim, a conclusão se sustenta. Fixado qualquer \(\tau\),
apenas o terceiro termo do \(\min\) que define \(m_i(\tau;N)\) depende
de \(N\), e depende de forma crescente (\(m_i^1+\mu_i^-(N-1-\tau)\)
cresce com \(N\) quando \(\mu_i^->0\)); os outros dois termos ficam
parados. Um mínimo em que um termo só sobe e os demais são fixos não
pode diminuir, logo \(m_i(\tau;N)\) é não-decrescente em \(N\) para cada
\(\tau\) fixo --- as semanas ``antigas'' podem crescer ao alongar o
horizonte, nunca encolher. Como \(\eta_{s,i\tau}\) é um max-fluxo com a
capacidade de fonte escalada por \(m_i(\tau;N)\) (§I.3.2), e max-fluxo é
não-decrescente em qualquer capacidade de aresta,
\(\varepsilon_{i\tau}(N)\) herda essa monotonicidade, e portanto também
cada \(q_{i\tau}(N)=F_i(m_i(\tau;N))\varepsilon_{i\tau}^{exog}(N)h_id_i\).
Somando as \(N\) semanas ``antigas'' (que só podem ter crescido) à
semana nova --- que cai exatamente no piso \(m_i^1\ge1\) e contribui
\(\rho_i\varepsilon_{i,N}h_id_im_i^1>0\), salvo calendário morto ---,
obtém-se \(\Pi(N+1)>\Pi(N)\): \(\Pi\) é estritamente crescente a cada
passo, inclusive no caso degenerado retangular, onde a semana nova
sozinha já garante o aumento. Uma função inteira estritamente monotônica
cruza qualquer limiar \(Q_i\) em exatamente um ponto, o que garante a
convergência da busca por incremento simples (Parte XIV, Fase 1) --- a
busca binária é uma alternativa equivalente quando o custo de
reavaliação por semana for relevante.

O perfil testado em cada tentativa não precisa estar ``certo'' antes de
\(N_i\) convergir: ele é, por definição, o perfil correto para aquele
\(N\) testado, e é descartado assim que falha o teste de limiar --- nada
é reaproveitado entre tentativas. O perfil que sobra, construído em
\(N=N_i^\star\), é o único que os cálculos das partes seguintes
utilizam: \(\{m_{ij}\} \leftarrow \text{RampProfile.build}(N_i, W_i)\) é
exato, sem resíduo das tentativas rejeitadas. Restrição de consistência:
\(r + d \le N_i\).

Um detalhe de ordem é necessário para que essa busca sequer possa
consultar \(\eta_{p,j}, \eta_{w,j}\) (indexados pela semana
\textbf{absoluta} \(j\)): \(j_i^0\) precisa estar definido antes de a
busca por \(N_i\) começar. Como \(j_i^0\) (§I.5) depende apenas de
atividades predecessoras --- já resolvidas pela ordem topológica em que
as atividades são processadas ---, nunca de \(N_i\) ou \(\{m_{ij}\}\) da
própria atividade \(i\), ele pode e deve ser calculado primeiro, sem
gerar dependência circular alguma com a busca que vem a seguir (Parte
XIV, Fase 1).

\hypertarget{i.5-predecessores-e-completude-parcial}{%
\subsubsection{I.5 Predecessores e completude
parcial}\label{i.5-predecessores-e-completude-parcial}}

\[B_i \subseteq \mathcal{A}\setminus\{A_i\}, \qquad \theta_i : B_i \to (0,1]\]

Para \(b\in B_i\), \(\theta_i(b)\) é a fração de escopo de \(b\) exigida
antes que \(A_i\) possa iniciar.

\textbf{Medida de completude em quantitativo:}

\[\Phi_b(j) = \frac{1}{Q_b}\sum_{j'=j_b^0}^{j} q_{b,j'}\]

\textbf{Restrição de início:}

\[\boxed{j_i^0 \;\ge\; \max_{b\in B_i}\ \min\!\{j : \Phi_b(j) \ge \theta_i(b)\}}\]

\textbf{Tipo de \(\theta_i(b)\)} --- declarar por aresta, pois a
semântica difere sob \(\varepsilon_{ij}\) variável:

\begin{longtable}[]{@{}
  >{\raggedright\arraybackslash}p{(\columnwidth - 6\tabcolsep) * \real{0.2500}}
  >{\raggedright\arraybackslash}p{(\columnwidth - 6\tabcolsep) * \real{0.2500}}
  >{\raggedright\arraybackslash}p{(\columnwidth - 6\tabcolsep) * \real{0.2500}}
  >{\raggedright\arraybackslash}p{(\columnwidth - 6\tabcolsep) * \real{0.2500}}@{}}
\toprule\noalign{}
\begin{minipage}[b]{\linewidth}\raggedright
Tipo
\end{minipage} & \begin{minipage}[b]{\linewidth}\raggedright
Condição de início
\end{minipage} & \begin{minipage}[b]{\linewidth}\raggedright
Formulação
\end{minipage} & \begin{minipage}[b]{\linewidth}\raggedright
Quando usar
\end{minipage} \\
\midrule\noalign{}
\endhead
\bottomrule\noalign{}
\endlastfoot
\textbf{Q} (quantidade) &
\(\sum_{j'\le j} q_{b,j'} \ge \theta_i(b)\,Q_b\) &
\(\Phi_b(j) \ge \theta_i(b)\) & progresso físico real (defensável) \\
\textbf{T} (tempo) & \(j \ge j_b^0 + \theta_i(b)\,N_b\) & fração de
calendário & representação Gantt / MS Project \\
\end{longtable}

Sob \(\varepsilon_{ij}\) variável as duas \textbf{não coincidem}: tipo T
libera o sucessor mais cedo que o fisicamente entregue quando há semanas
improdutivas. Tipo Q é o adotado neste formulário --- restrição exata
via rastreamento da produção acumulada realizada \(q_{b,j'}\) no solver,
incluindo o fechamento fracionário da última semana (Parte IV). O caso
\(\theta_i(b)=1\) recupera a precedência fim-início clássica em ambos os
tipos.

\begin{center}\rule{0.5\linewidth}{0.5pt}\end{center}

\hypertarget{parte-ii-paruxe2metros-e-constantes-globais}{%
\subsection{Parte II --- Parâmetros e Constantes
Globais}\label{parte-ii-paruxe2metros-e-constantes-globais}}

Esta seção reúne constantes e faixas de valores compartilhadas por todo
o projeto. Nenhuma delas é calculada a partir de outra parte do
formulário --- todas são definidas pelo usuário ou pela política
comercial da empresa antes de o cálculo começar.

\[s_{\text{mes}} = \frac{365\times3+366}{4\times12\times7} = \frac{1461}{336} \approx 4{,}3482 \ \text{semanas/mês}\]

\(s_{\text{mes}}\) converte uma taxa mensal em uma taxa semanal. Como
nem todo mês tem o mesmo número de semanas, e o ano tem, em média, um
pouco mais que 52 semanas por causa dos anos bissextos, usar
simplesmente ``4 semanas por mês'' introduziria um viés sistemático ao
longo do tempo. Esta constante usa a média de um ciclo de quatro anos
--- três anos comuns de 365 dias mais um bissexto de 366 --- para obter
o fator de conversão correto. É o mesmo número usado, por exemplo, para
diluir o aluguel mensal de um equipamento (§I.3.1) em um custo por
semana.

\[d_i \in \{4,5,6\}\ \text{dias úteis/semana}, \qquad h_i \in \{6,8\}\ \text{horas/dia}, \qquad h_id_i \equiv \text{horas/semana da equipe }i\]

\(d_i\) e \(h_i\) descrevem o regime de jornada de cada equipe: quantos
dias por semana e quantas horas por dia ela efetivamente opera. Esses
valores variam por atividade porque dependem de acordos trabalhistas,
turnos e da política operacional de cada frente --- uma equipe em regime
de 6 dias/semana com turnos de 8 horas tem uma capacidade semanal bem
diferente de uma em 4 dias/semana com turnos de 6 horas, mesmo que o
custo mensal contratual seja parecido.

Taxas de BDI (entrada; valores de referência do modelo AGIS entre
parênteses):

\[\beta^{EQ},\ \beta^{MO}\ (\approx29{,}97\%),\ \beta^{MAT},\ \beta^{COMB}\ (\approx4\%),\ \beta^{IND}\ (\approx25{,}21\%)\]

Os \(\beta\) são as taxas de BDI (Benefícios e Despesas Indiretas)
associadas a cada categoria de custo: o percentual que se soma ao custo
direto para cobrir despesas administrativas, riscos e lucro. Cada
categoria carrega uma taxa diferente porque a estrutura de custo
indireto por trás dela é diferente --- mão de obra, por exemplo,
tradicionalmente carrega encargos trabalhistas e riscos de gestão de
pessoal mais pesados que a simples compra de um material bruto, o que
explica \(\beta^{MO}\) ser muito maior que \(\beta^{MAT}\) nos valores
de referência.

\begin{center}\rule{0.5\linewidth}{0.5pt}\end{center}

\hypertarget{parte-iii-custo-horuxe1rio-de-catuxe1logo-gamma}{%
\subsection{\texorpdfstring{Parte III --- Custo Horário de Catálogo
(\(\gamma\))}{Parte III --- Custo Horário de Catálogo (\textbackslash gamma)}}\label{parte-iii-custo-horuxe1rio-de-catuxe1logo-gamma}}

O custo de catálogo é o preço ``de tabela'' de uma equipe: quanto ela
custaria por hora se operasse exatamente como planejado, sem atrasos,
perdas de eficiência ou imprevistos de canteiro. É o ponto de partida
antes de qualquer ajuste pela realidade da obra --- as próximas partes
mostram como esse valor se transforma quando a produção real entra em
cena.

\[\gamma_i^{EQ} = \sum_{r\in R_i,\,\text{cat}(r)=EQ} c_r n_r \qquad \text{(análogo para } MO, MAT, COMB\text{)}\]

Somando o custo horário de cada recurso (\(c_r\), calculado em §I.3.1)
multiplicado por quantos desse recurso existem em uma frente (\(n_r\)),
chega-se ao custo horário de cada categoria.

\[\gamma_i^{\mathrm{K,cat}} \equiv \gamma_i^{EQ}+\gamma_i^{MO} \qquad\text{[contratual: cobrado por semana mobilizada]}\]
\[\gamma_i^{\mathrm{V,cat}} \equiv \gamma_i^{MAT}+\gamma_i^{COMB} \qquad\text{[consumível: escala com produção]}\]
\[\gamma_i^{\mathrm{cat}} = \gamma_i^{\mathrm{K,cat}}+\gamma_i^{\mathrm{V,cat}}\qquad\text{[R\$/h]}\]

Agrupando as quatro categorias nos dois blocos de §0, chega-se ao custo
horário total ``de tabela'' da equipe --- e à divisão entre a parte que
é fixa por semana (K) e a parte que varia com a produção (V).

\[f_i^{\mathrm V} = \gamma_i^{\mathrm{V,cat}}/\gamma_i^{\mathrm{cat}} \qquad\text{[fração consumível do custo de catálogo]}\]

\(f_i^{\mathrm V}\) resume, num único número, o ``perfil de risco'' da
equipe: uma equipe com \(f_i^{\mathrm V}\) alto é mais sensível ao preço
de combustível e material; uma equipe com \(f_i^{\mathrm V}\) baixo é
mais sensível a quanto tempo permanece mobilizada, independentemente do
quanto produz.

\[p_i^{\mathrm{cat}} = \gamma_i^{\mathrm{cat}}/\rho_i \qquad\text{[R\$/un — preço unitário de catálogo]}\]

\[\boxed{C_i^{\mathrm{cat}} \equiv p_i^{\mathrm{cat}}\,Q_i = \frac{\gamma_i^{\mathrm{cat}}\,Q_i}{\rho_i}}\qquad\text{[R\$]}\]

\(C_i^{\mathrm{cat}}\) é o custo de catálogo da atividade --- base para
diagnósticos e para o envelope de catálogo do projeto inteiro (Parte
XIII). Dividindo o custo horário pela produção horária (\(\rho_i\),
§I.3.2), chega-se ao preço de tabela por unidade produzida;
multiplicando pelo quantitativo total da atividade, chega-se ao custo de
catálogo da atividade inteira --- o número contra o qual toda
ineficiência real será depois comparada.

\begin{center}\rule{0.5\linewidth}{0.5pt}\end{center}

\hypertarget{parte-iv-cronograma-semanal}{%
\subsection{Parte IV --- Cronograma
Semanal}\label{parte-iv-cronograma-semanal}}

Até aqui, os custos foram tratados como se toda semana fosse igualmente
produtiva. Esta seção introduz o calendário real: quantas unidades cada
atividade efetivamente produz, semana a semana, considerando quantas
frentes estão mobilizadas e quão eficiente é o trabalho naquela semana
específica.

Capacidade física e fatores exógenos semanais:

\[\boxed{B_{ij}=F_i(m_{ij})h_id_i\eta_{p,j}\eta_{w,j},
\quad q_{ij}=\min\{(Q_i-\sum_{k<j}q_{ik})_+,B_{ij}\}}\]

Não aplicar nenhuma penalização adicional por frentes não mobilizadas.

\textbf{Fechamento exato de escopo.} Seja \(j_i^1 \equiv j_i^0+N_i-1\) a
última semana ativa de \(A_i\), obtida pela busca monotônica em \(N_i\)
(Parte XIV, Fase 1) sobre o perfil nominal \(\{m_i(\tau)\}\) de \(W_i\).
Nas semanas interiores \(j<j_i^1\), a produção usa o perfil nominal
integralmente. Na última semana, o quantitativo residual fecha
exatamente o escopo:

\[q_{i,j_i^1} = Q_i - \sum_{j=j_i^0}^{j_i^1-1} q_{ij}\]

Por construção da busca monotônica (\(N_i\) é o menor inteiro tal que a
produção acumulada nominal atinge \(Q_i\), e a semana adicionada mais
recentemente é a primeira em que isso ocorre), o residual satisfaz
\(0 < q_{i,j_i^1} \le \rho_i\varepsilon_{i,j_i^1}h_id_im_i(N_i-1)\) ---
nunca excede a produção nominal da última semana, e a semana de
fechamento sempre tem \(\varepsilon_{i,j_i^1}>0\) (uma semana com
\(\varepsilon=0\) não altera a produção acumulada, de modo que a busca
automaticamente avança para a semana seguinte). Define-se a
\textbf{fração de fechamento}:

\[\varphi_i \equiv \frac{q_{i,j_i^1}}{F_i(m_i(N_i-1))\,h_id_i\,\eta_{p,j_i^1}\eta_{w,j_i^1}} \in (0,1]\]

representando a fração da capacidade nominal da última semana
efetivamente utilizada para completar o escopo. A frota mobilizada
(\(m_{ij}\), e portanto o custo contratual \(C_{ij}^{\mathrm{K,op}}\),
Parte V) permanece no valor nominal do perfil na última semana ---
recursos contratuais (EQ/MO) são cobrados pela semana mobilizada
independentemente da fração de produção efetivamente necessária naquela
semana (§0). Apenas o consumível é prorateado pela produção realizada,
via \(\varphi_i\), garantindo o fecho exato do escopo:

\[\sum_{j=j_i^0}^{j_i^1} q_{ij} = Q_i \qquad\text{(exato, por construção)}\]

Contadores: \(M_i \equiv \sum_j m_{ij}\) (frente-semanas totais, ao
valor nominal --- inclui a última semana integralmente),
\(N_i = |\{j:m_{ij}>0\}|\) (semanas de calendário ativas).

\[\mathcal{A}_j \coloneqq \{A_i\in\mathcal{A} : j\in[j_i^0,\,j_i^0+N_i-1]\}\]

\begin{center}\rule{0.5\linewidth}{0.5pt}\end{center}

\hypertarget{parte-v-custo-direto-semanal-sem-bdi}{%
\subsection{Parte V --- Custo Direto Semanal (sem
BDI)}\label{parte-v-custo-direto-semanal-sem-bdi}}

Esta seção converte os quantitativos semanais da Parte IV em valores
monetários, separando o que é cobrado independentemente da produção (o
bloco contratual, K) do que acompanha exatamente o que foi produzido (o
bloco consumível, V) --- a mesma distinção introduzida em §0.

\[C_{ij}^{\mathrm{K,op}} = \gamma_i^{\mathrm{K,cat}}\,h_id_i\,m_{ij}\qquad\text{[R\$] — nominal em toda semana ativa, inclusive a última}\]

\(C_{ij}^{\mathrm{K,op}}\) é o custo contratual da semana: o valor pago
pelas frentes mobilizadas, que depende apenas do regime de jornada e de
quantas frentes estão em campo --- não de quanto foi produzido. É como
alugar um equipamento por semana: o valor é o mesmo, tenha a equipe
produzido muito ou pouco naquela semana.

\[C_{ij}^{\mathrm{V,op}} = \gamma_i^{\mathrm{V,cat}}\cdot\frac{q_{ij}}{\rho_i}\qquad\text{[R\$]}\]

\(C_{ij}^{\mathrm{V,op}}\) é o custo consumível da semana: proporcional
ao quanto foi de fato produzido (\(q_{ij}\)), e por isso já incorpora,
de forma implícita, tanto a eficiência da semana quanto o fechamento
fracionário da última semana (Parte IV). É como pagar por combustível:
quanto mais se produz, mais se consome, e vice-versa.

Nas semanas interiores,
\(q_{ij}/\rho_i \equiv \varepsilon_{ij}h_id_im_{ij}\) por definição de
\(q_{ij}\) (Parte IV) --- a formulação acima generaliza a forma
tempo-a-tempo e permanece exata na semana de fechamento, onde \(q_{ij}\)
é o residual e não \(\varepsilon_{ij}h_id_im_{ij}\) diretamente.

\[C_{ij}^{\mathrm{op}} = C_{ij}^{\mathrm{K,op}}+C_{ij}^{\mathrm{V,op}}\qquad\text{[R\$]}\]

Somando os dois blocos, chega-se ao custo direto daquela semana
específica --- a peça básica a partir da qual todos os totais das
próximas partes são construídos.

\begin{center}\rule{0.5\linewidth}{0.5pt}\end{center}

\hypertarget{parte-vi-custo-direto-total-por-atividade-sem-bdi}{%
\subsection{Parte VI --- Custo Direto Total por Atividade (sem
BDI)}\label{parte-vi-custo-direto-total-por-atividade-sem-bdi}}

Somando os custos semanais da Parte V ao longo de toda a duração da
atividade, chega-se ao seu custo direto total --- e a uma das
identidades mais úteis deste formulário, detalhada logo abaixo.

\[C_i^{\mathrm{K,cat}} = \frac{Q_i\,\gamma_i^{\mathrm{K,cat}}}{\rho_i}\qquad\text{[R\$] — catálogo, parcela contratual, à eficiência nominal}\]
\[C_i^{\mathrm{V,cat}} = \frac{Q_i\,\gamma_i^{\mathrm{V,cat}}}{\rho_i}\qquad\text{[R\$] — catálogo, parcela consumível}\]
\[C_i^{\mathrm{K,op}} = \sum_{j} C_{ij}^{\mathrm{K,op}} = \gamma_i^{\mathrm{K,cat}}\,h_id_i\,M_i\qquad\text{[R\$]}\]
\[C_i^{\mathrm{V,op}} = \sum_{j} C_{ij}^{\mathrm{V,op}} = \frac{\gamma_i^{\mathrm{V,cat}}}{\rho_i}\sum_j q_{ij}\qquad\text{[R\$]}\]
\[C_i^{\mathrm{op}} = C_i^{\mathrm{K,op}}+C_i^{\mathrm{V,op}}\qquad\text{[R\$]}\]
\[p_i^{\mathrm{op}} = C_i^{\mathrm{op}}/Q_i\qquad\text{[R\$/un]}\]

\textbf{Identidade --- invariância do consumível.} Usando o fecho de
escopo exato \(\sum_jq_{ij}=Q_i\) (Parte IV):

\[\boxed{C_i^{\mathrm{V,op}} \equiv C_i^{\mathrm{V,cat}} = \frac{Q_i\,\gamma_i^{\mathrm{V,cat}}}{\rho_i}}\qquad\text{(exato, por construção do fechamento de escopo)}\]

O consumível operacional \textbf{é} o consumível de catálogo --- não há
inflação por ineficiência. O total de material/combustível depende
apenas de \(Q_i\), nunca de quantas semanas foram necessárias para
executar. Esta identidade \textbf{não vale} para o contratual:
\(C_i^{\mathrm{K,op}}=\gamma_i^{\mathrm{K,cat}}h_id_iM_i\) cresce quando
a eficiência piora (mais frente-semanas \(M_i\) para fechar o mesmo
\(Q_i\), já que \(M_i\) conta a semana de fechamento integralmente, não
prorateada).

\textbf{Desvio de catálogo por atividade.} Comparando
\(C_i^{\mathrm{op}}\) com \(C_i^{\mathrm{cat}}\) (Parte III), obtém-se
uma leitura de quanto a atividade se desviou do plano de catálogo --- em
termos relativos e absolutos:

\[R_i \equiv \frac{C_i^{\mathrm{op}}}{C_i^{\mathrm{cat}}} \qquad\qquad \boxed{\Delta_i^{\mathrm{cat}} \equiv C_i^{\mathrm{op}}-C_i^{\mathrm{cat}} = (R_i-1)\,C_i^{\mathrm{cat}}}\qquad\text{[R\$]}\]

\(\Delta_i^{\mathrm{cat}}>0\) significa que a atividade custou mais do
que o catálogo previa; \(\Delta_i^{\mathrm{cat}}<0\), que custou menos.
\(R_i\) é a mesma leitura em forma adimensional, mais direta para
comparar o grau de desvio entre atividades de portes diferentes --- é o
análogo, em nível de atividade, do
\(R=C^{\mathrm{tot}}/C^{\mathrm{cat}}\) de projeto (Parte XIII). Sua
forma fechada, em função da eficiência média realizada, é dada na Parte
VIII, logo após \(\bar\varepsilon_i\) ser definida com exatidão.

Uma advertência sobre como (não) recalcular esse desvio sob uma
produtividade revisada --- por exemplo, ao incorporar um \(\rho_i\)
observado em campo, diferente do calibrado originalmente: a identidade
\(C_i^{\mathrm{V,op}}\equiv\gamma_i^{\mathrm{V,cat}}Q_i/\rho_i\) (acima)
é, por construção, invariante ao cronograma, então substituir \(\rho_i\)
por um \(\rho_i'\) revisado nessa fórmula sozinha ainda produz um número
correto. Mas \(C_i^{\mathrm{K,op}}=\gamma_i^{\mathrm{K,cat}}h_id_iM_i\)
depende de \(M_i\), resultado da busca por \(N_i\) (Parte XIV, Fase 1)
\emph{sob o \(\rho_i\) original}. Trocar \(\rho_i\) por \(\rho_i'\) nas
fórmulas de custo sem refazer essa busca deixa \(M_i\) preso a um
cronograma que, sob a nova produtividade, não fecha mais o escopo
\(Q_i\) (Parte IV) --- o modelo estaria implicitamente descrevendo uma
atividade que produz uma quantidade diferente de \(Q_i\). Recalcular
\(\Delta_i^{\mathrm{cat}}\) sob produtividade revisada exige rederivar
\(N_i\), \(\{m_{ij}\}\) e \(\varphi_i\) a partir de \(\rho_i'\) (Fases
0--3), não apenas substituir \(\rho_i\) nas fórmulas de custo mantendo o
cronograma antigo.

\begin{center}\rule{0.5\linewidth}{0.5pt}\end{center}

\hypertarget{parte-vii-razuxf5es-operacionalcatuxe1logo-semanais}{%
\subsection{Parte VII --- Razões Operacional/Catálogo
(semanais)}\label{parte-vii-razuxf5es-operacionalcatuxe1logo-semanais}}

As razões desta seção comparam o que a atividade realmente custou e
produziu em uma dada semana com o que o catálogo previa --- funcionam
como um termômetro semanal de quão perto (ou longe) a operação ficou do
planejado.

Custo semanal de catálogo (pool e consumível à taxa nominal, sem
desconto de \(\varepsilon_{ij}\)):

\[C_{ij}^{\mathrm{cat}} = \gamma_i^{\mathrm{cat}}\,h_id_i\,m_{ij}\qquad\text{[R\$]}\]

Preço unitário operacional semanal:

\[p_{ij}^{\mathrm{op}} = C_{ij}^{\mathrm{op}}/q_{ij}\qquad\text{[R\$/un]}\]

Razões:

\[\zeta_{ij} = \frac{C_{ij}^{\mathrm{op}}}{C_{ij}^{\mathrm{cat}}} = \frac{\gamma_i^{\mathrm{K,cat}}+\varepsilon_{ij}\gamma_i^{\mathrm{V,cat}}}{\gamma_i^{\mathrm{cat}}}\qquad[\ ]\]

\[\pi_{ij} = \frac{p_{ij}^{\mathrm{op}}}{p_i^{\mathrm{cat}}} = \frac{\gamma_i^{\mathrm{K,cat}}+\varepsilon_{ij}\gamma_i^{\mathrm{V,cat}}}{\varepsilon_{ij}\,\gamma_i^{\mathrm{cat}}}\qquad[\ ]\]

\[\boxed{\zeta_{ij} = \varepsilon_{ij}\cdot\pi_{ij}}\]

Em palavras simples: \(\zeta_{ij}\) compara custos totais --- quanto se
gastou frente ao previsto ---, enquanto \(\pi_{ij}\) compara preços
unitários --- quanto custou cada unidade produzida frente ao preço de
tabela. Como uma semana ineficiente ainda paga o custo contratual cheio
mas produz menos, \(\pi_{ij}\) tende a subir quando \(\varepsilon_{ij}\)
cai: produzir menos encarece cada unidade entregue.

Usando a fração consumível \(f_i^{\mathrm V}\) (Parte III):

\[\zeta_{ij} = 1-(1-\varepsilon_{ij})\,f_i^{\mathrm V}, \qquad \pi_{ij} = \frac{1}{\varepsilon_{ij}}-\left(\frac{1}{\varepsilon_{ij}}-1\right)f_i^{\mathrm V}\]

\(\zeta_{ij}\) é diagnóstico local por semana --- não deve ser somado
ingenuamente para recuperar um ``\(\zeta_i\)'' de atividade.

\begin{center}\rule{0.5\linewidth}{0.5pt}\end{center}

\hypertarget{parte-viii-camada-bdi-direta-por-atividade}{%
\subsection{Parte VIII --- Camada BDI (direta, por
atividade)}\label{parte-viii-camada-bdi-direta-por-atividade}}

Esta seção aplica as taxas de BDI da Parte II sobre os custos de
catálogo e operacionais já calculados, chegando ao custo de cada
atividade como ele efetivamente entra na proposta comercial --- e a um
BDI equivalente por atividade, que varia semana a semana conforme muda a
mistura entre custo contratual e consumível.

\[\gamma_i^{\mathrm{K,cat},\beta} = (1+\beta^{\mathrm{EQ}})\gamma_i^{\mathrm{EQ}}+(1+\beta^{\mathrm{MO}})\gamma_i^{\mathrm{MO}}\]
\[\gamma_i^{\mathrm{V,cat},\beta} = (1+\beta^{\mathrm{MAT}})\gamma_i^{\mathrm{MAT}}+(1+\beta^{\mathrm{COMB}})\gamma_i^{\mathrm{COMB}}\]
\[\gamma_i^{\mathrm{cat},\beta} = \gamma_i^{\mathrm{K,cat},\beta}+\gamma_i^{\mathrm{V,cat},\beta}\]

\[C_{ij}^{\mathrm{K,op},\beta} = \gamma_i^{\mathrm{K,cat},\beta}\,h_id_i\,m_{ij}\]
\[C_{ij}^{\mathrm{V,op},\beta} = \gamma_i^{\mathrm{V,cat},\beta}\cdot\frac{q_{ij}}{\rho_i}\]
\[C_{ij}^{\mathrm{op},\beta} = C_{ij}^{\mathrm{K,op},\beta}+C_{ij}^{\mathrm{V,op},\beta}\]

\textbf{Taxas ponderadas por pool.} O BDI equivalente de cada bloco (K e
V) é a média de \(\beta\) das suas categorias, ponderada pela composição
de custo de catálogo da equipe:

\[\beta_i^{\mathrm K} \equiv \frac{\gamma_i^{\mathrm{EQ}}\beta^{\mathrm{EQ}}+\gamma_i^{\mathrm{MO}}\beta^{\mathrm{MO}}}{\gamma_i^{\mathrm{K,cat}}}, \qquad \beta_i^{\mathrm V} \equiv \frac{\gamma_i^{\mathrm{MAT}}\beta^{\mathrm{MAT}}+\gamma_i^{\mathrm{COMB}}\beta^{\mathrm{COMB}}}{\gamma_i^{\mathrm{V,cat}}}\]

de modo que, por construção,
\(\gamma_i^{\mathrm{K,cat},\beta}=\gamma_i^{\mathrm{K,cat}}(1+\beta_i^{\mathrm K})\)
e
\(\gamma_i^{\mathrm{V,cat},\beta}=\gamma_i^{\mathrm{V,cat}}(1+\beta_i^{\mathrm V})\).

Fração consumível efetiva da semana \(j\):

\[\delta_{ij} \equiv \frac{\varepsilon_{ij}\,\gamma_i^{\mathrm{V,cat}}}{\gamma_i^{\mathrm{K,cat}}+\varepsilon_{ij}\,\gamma_i^{\mathrm{V,cat}}}\]

\[\boxed{\beta_{ij} = (1-\delta_{ij})\,\beta_i^{\mathrm K} + \delta_{ij}\,\beta_i^{\mathrm V}}\]

Decrescente em \(\varepsilon_{ij}\) sempre que
\(\beta_i^{\mathrm V}<\beta_i^{\mathrm K}\) (caso típico, dado
\(\beta^{MAT},\beta^{COMB}\ll\beta^{EQ},\beta^{MO}\) nos valores de
referência): maior eficiência aloca mais custo ao bloco consumível, de
BDI mais baixo, reduzindo o BDI equivalente semanal.

\textbf{Totais por atividade:}

\[C_i^{\mathrm{K,op},\beta} = \gamma_i^{\mathrm{K,cat},\beta}\,h_id_i\,M_i \qquad C_i^{\mathrm{V,op},\beta} = \gamma_i^{\mathrm{V,cat},\beta}\cdot\frac{Q_i}{\rho_i}\]

\[C_i^{D,\beta} \equiv C_i^{\mathrm{K,op},\beta}+C_i^{\mathrm{V,op},\beta} \qquad\text{[R\$] — total direto da atividade com BDI}\]

\textbf{Eficiência média ponderada por frente-semana} --- agrega o
efeito de \(\varepsilon_{ij}\) ao longo de toda a atividade em um único
fator. Definida diretamente a partir do fecho de escopo exato (Parte
IV), a agregação vale mesmo com o fechamento fracionário da última
semana:

\[\bar\varepsilon_i \equiv \frac{Q_i}{\rho_i\,h_id_i\,M_i}\]

Fora da semana de fechamento, isso equivale à soma
\(\sum_j\varepsilon_{ij}m_{ij}/M_i\) termo a termo; na semana de
fechamento, o termo correspondente entra como
\(\varphi_i\varepsilon_{i,j_i^1}m_{i,j_i^1}\) em vez de
\(\varepsilon_{i,j_i^1}m_{i,j_i^1}\) --- a fração de fechamento
\(\varphi_i\) (Parte IV) é exatamente essa correção, de modo que a
definição acima já a incorpora, sem tratamento à parte.

\[\delta_i = \frac{\bar\varepsilon_i\,\gamma_i^{\mathrm{V,cat}}}{\gamma_i^{\mathrm{K,cat}}+\bar\varepsilon_i\,\gamma_i^{\mathrm{V,cat}}} \;\left(=\frac{C_i^{\mathrm{V,op}}}{C_i^{\mathrm{op}}}\text{, exato}\right)\]

\[\boxed{\beta_i^{\mathrm{eq}} \equiv \frac{C_i^{D,\beta}}{C_i^{\mathrm{op}}}-1 = (1-\delta_i)\,\beta_i^{\mathrm K}+\delta_i\,\beta_i^{\mathrm V}}\]

Com \(\bar\varepsilon_i\) exato, o desvio de catálogo por atividade
\(R_i\equiv C_i^{\mathrm{op}}/C_i^{\mathrm{cat}}\) (Parte VI) também
ganha forma fechada:

\[\boxed{R_i = f_i^{\mathrm V} + \frac{1-f_i^{\mathrm V}}{\bar\varepsilon_i}}\]

\(R_i=1\) exatamente quando \(\bar\varepsilon_i=1\): todo o desvio de
catálogo nasce da diferença entre a eficiência média realizada e a
eficiência nominal de catálogo --- nunca da composição K/V da equipe
isoladamente.

\begin{center}\rule{0.5\linewidth}{0.5pt}\end{center}

\hypertarget{parte-ix-pool-de-recursos-contratuais-perpetuauxe7uxe3o}{%
\subsection{Parte IX --- Pool de Recursos Contratuais
(Perpetuação)}\label{parte-ix-pool-de-recursos-contratuais-perpetuauxe7uxe3o}}

Até aqui, cada atividade foi tratada isoladamente. Na prática, porém, o
mesmo equipamento ou a mesma equipe frequentemente passa de uma
atividade para outra sem ser desmobilizado, evitando o custo e o tempo
de tirá-lo do canteiro e trazê-lo de volta depois. Esta seção trata
desse aproveitamento entre atividades: quando vale a pena manter um
recurso parado por alguns dias esperando a próxima frente de trabalho, e
quando é mais barato desmobilizar e remobilizar.

\hypertarget{ix.1-tipos-instuxe2ncias-e-demanda-semanal}{%
\subsubsection{IX.1 Tipos, instâncias e demanda
semanal}\label{ix.1-tipos-instuxe2ncias-e-demanda-semanal}}

\(\mathcal{T}\) = conjunto de tipos de recurso (união de categorias e
especificação sobre todos os \(R_i\)).

\textbf{Demanda semanal por tipo:}

\[n_j(\tau) = \sum_{i\,:\,j\in[j_i^0,\,j_i^0+N_i-1]} m_{ij}\cdot n_{\tau}^{i}\]

\textbf{Pool físico} \(\{r\}\): conjunto de instâncias com
\(\text{type}(r)\in\mathcal{T}\), taxa semanal \(c_r^{\text{sem}}\)
(igual em semanas ativas e ociosas --- §IX.3), e custos
\(c_r^{\text{mob}}\), \(c_r^{\text{demob}}\).

\[qt_\tau = \max_j n_j(\tau) \qquad\text{[tamanho do pool de tipo }\tau\text{]}\]

\hypertarget{ix.2-problema-de-minimizauxe7uxe3o-e-condiuxe7uxe3o-de-perpetuauxe7uxe3o}{%
\subsubsection{IX.2 Problema de minimização e condição de
perpetuação}\label{ix.2-problema-de-minimizauxe7uxe3o-e-condiuxe7uxe3o-de-perpetuauxe7uxe3o}}

\[\min_{\{r\},\{[j_r^0,j_r^1]\}} \ \underbrace{\sum_r\sum_{j=j_r^0}^{j_r^1} c_r^{\text{ativo/standby}}(j)}_{\text{permanência}} + \underbrace{\sum_r (c_r^{\text{mob}}+c_r^{\text{demob}})}_{\text{mob/demob}}\]

\[\text{s.a.}\quad \bigl|\{r:\text{type}(r)=\tau,\ j\in[j_r^0,j_r^1]\}\bigr| \ge n_j(\tau) \quad\forall j,\,\forall\tau\]

Para tipo único fungível (instâncias perfeitamente substituíveis, com
custo de permanência e de mob/demob homogêneos), a Parte XIV (Fase 4)
apresenta um algoritmo exato por decomposição em níveis do perfil de
demanda.

\textbf{Condição de perpetuação} (entre duas alocações consecutivas, com
gap em semanas):

\[\text{gap}_{uv} < G_r^* \equiv \frac{c_r^{\text{mob}}+c_r^{\text{demob}}}{c_r^{\text{sem}}} \quad\Longleftrightarrow\quad c_r^{\text{mob}}+c_r^{\text{demob}} > c_r^{\text{sem}}\cdot\text{gap}_{uv}\qquad\text{[$c_r^{\text{sem}}$ em R\$/sem]}\]

\textbf{Cap absoluto:}
\(\text{gap}_{uv} > G_r^{\text{viável}} \Rightarrow\) não perpetuar,
independente de custo.

\textbf{Estrutura de cadeia}
\(\mathcal{C}_p = (i_1 \to i_2 \to \cdots \to i_{n_p})\): cada par
consecutivo \((i_k, i_{k+1})\) satisfaz o critério. Uma cadeia gera
\textbf{um} par \((c_r^{\text{mob}}, c_r^{\text{demob}})\) mais standby
nos gaps internos.

\hypertarget{ix.3-frauxe7uxf5es-invariantes-por-recurso}{%
\subsubsection{IX.3 Frações invariantes por
recurso}\label{ix.3-frauxe7uxf5es-invariantes-por-recurso}}

Como o custo dos recursos contratuais é definido a partir do
\textbf{tempo de permanência} (\(c_r^{\text{sem}}\) por semana ativa,
\(c_r^{\text{standby}}\) por semana de gap), a separação em fração ativa
e fração standby é um \textbf{invariante temporal} de cada instância ---
completamente determinado pela cadeia de perpetuação, sem dependência de
parâmetros de custo.

Para cada instância \(r\) com cadeia
\(\mathcal{C}_p = (i_1 \to \cdots \to i_{n_p})\):

\[T_r^{\text{ativo}} = \sum_{k=1}^{n_p} N_{i_k} \qquad\text{[semanas produtivas]}\]

\[T_r^{\text{standby}} = \sum_{k=1}^{n_p-1} \text{gap}_{i_k,i_{k+1}} \qquad\text{[semanas de standby interno]}\]

\[T_r^{\text{perm}} = T_r^{\text{ativo}} + T_r^{\text{standby}} \qquad\text{[total de semanas de permanência]}\]

\[\boxed{f_r^{\text{standby}} = \frac{T_r^{\text{standby}}}{T_r^{\text{perm}}} \in [0,1), \qquad f_r^{\text{ativo}} = 1 - f_r^{\text{standby}}}\]

\textbf{Propriedade invariante:} dados o schedule \(\{(j_i^0, N_i)\}\) e
as cadeias de perpetuação, \(f_r^{\text{standby}}\) é completamente
determinado. Não depende de \(c_r^{\text{sem}}\),
\(c_r^{\text{standby}}\), \(c_r^{\text{mob}}\) nem
\(c_r^{\text{demob}}\) --- é uma grandeza puramente temporal, intrínseca
ao padrão de implantação do recurso.

\textbf{Casos:}

\begin{longtable}[]{@{}
  >{\raggedright\arraybackslash}p{(\columnwidth - 2\tabcolsep) * \real{0.5000}}
  >{\raggedright\arraybackslash}p{(\columnwidth - 2\tabcolsep) * \real{0.5000}}@{}}
\toprule\noalign{}
\begin{minipage}[b]{\linewidth}\raggedright
Situação
\end{minipage} & \begin{minipage}[b]{\linewidth}\raggedright
\(f_r^{\text{standby}}\)
\end{minipage} \\
\midrule\noalign{}
\endhead
\bottomrule\noalign{}
\endlastfoot
Cadeia singleton (\(n_p=1\), sem perpetuação) & \(0\) \\
Perpetuação com gaps pequenos & \(\in(0,1)\), próximo de 0 quando
atividades são densas \\
Recurso em standby quase todo tempo & \(\to 1\) (não deve ocorrer:
perpetuação só é racional quando gap \(< G_r^*\)) \\
\end{longtable}

\textbf{Taxa efetiva semanal.} A separação pool/estoque implica
\(c_r^{\text{standby}} = c_r^{\text{sem}}\): EQ (aluguel) e MO (salário
+ CLT) custam igual em semanas ativas e ociosas --- os consumíveis, que
escalariam com produção, já foram segregados em \(C^{\text{estoque}}\).
Portanto:

\[\boxed{\bar{c}_r = c_r^{\text{sem}} \qquad \text{(pool puro: standby = ativo)}}\]

A fração \(f_r^{\text{standby}}\) é \textbf{diagnóstico temporal} puro
--- não afeta o custo, quantifica o grau de ociosidade da instância. O
benefício da perpetuação é exclusivamente a redução de \(C^{MOB,eq,o}\)
(mob/demob evitados), ao custo de semanas adicionais de permanência a
\(c_r^{\text{sem}}\).

\hypertarget{ix.4-custo-do-pool}{%
\subsubsection{IX.4 Custo do pool}\label{ix.4-custo-do-pool}}

\[\boxed{C^{\text{pool},o} = \sum_r T_r^{\text{perm}} \cdot c_r^{\text{sem}}}\]

\[C^{\text{pool},\beta} = \sum_r (1+\beta_{\text{cat}(r)})\cdot T_r^{\text{perm}}\cdot c_r^{\text{sem}}\]

\textbf{Custo indireto de mob/demob do pool} (por cadeia, com
\(\beta^{IND}\)):

\[C^{MOB,eq,o} = \sum_r N_{\text{chains}}^r \cdot (c_r^{\text{mob}}+c_r^{\text{demob}}), \qquad C^{MOB,eq,\beta} = (1+\beta^{IND})\cdot C^{MOB,eq,o}\]

\textbf{Reconciliação com a Parte VI:}

\[\boxed{C^{\mathrm{pool}} \;\ge\; \sum_i C_i^{\mathrm{K,op}}}\]

com igualdade apenas se nenhum recurso é perpetuado entre atividades
distintas e \(f_r^{\text{standby}}=0\ \forall r\): a soma por atividade
conta apenas semanas ativas (\(M_i\)), enquanto
\(C^{\mathrm{pool}}=\sum_rT_r^{\text{perm}}c_r^{\text{sem}}\) inclui
também as semanas de standby interno das cadeias perpetuadas
(\(T_r^{\text{standby}}\)) --- custo real de permanência sem
contrapartida de produção. \(C_i^{\mathrm{K,op}}\) permanece grandeza de
diagnóstico por atividade (quanto cada atividade consumiria de recurso
contratual se operasse isoladamente); o desembolso contratual real do
projeto --- incluindo o custo de manter recursos parados entre
mobilizações --- é \(C^{\mathrm{pool}}\).

\textbf{Escopo do algoritmo de dimensionamento.} A decomposição por
níveis empregada na Parte XIV, Fase 4 é exata para um tipo de recurso
fungível homogêneo --- instâncias perfeitamente substituíveis entre si,
com custo de permanência e de mobilização/desmobilização idênticos.
Variantes multi-tipo com recursos não perfeitamente substituíveis
(capacidades ou custos heterogêneos dentro do mesmo tipo nominal)
constituem um problema de otimização combinatória distinto, de
complexidade superior à do caso fungível (cf.~Bertossi, Carraresi \&
Gallo 1987); nesses casos, os tipos devem ser subdivididos em subtipos
homogêneos antes de aplicar o algoritmo desta seção.

\begin{center}\rule{0.5\linewidth}{0.5pt}\end{center}

\hypertarget{parte-x-consumuxedveis-aquisiuxe7uxe3o-demanda-e-reposiuxe7uxe3o}{%
\subsection{Parte X --- Consumíveis: Aquisição, Demanda e
Reposição}\label{parte-x-consumuxedveis-aquisiuxe7uxe3o-demanda-e-reposiuxe7uxe3o}}

Consumíveis passam por três contas diferentes, ligadas por uma cadeia de
derivação: quanto custou comprar o material ou combustível efetivamente
usado na produção (aquisição, direto), quanto foi fisicamente consumido
semana a semana (demanda física, derivada do cronograma) e quanto custa
manter estoque parado no canteiro entre a compra e o consumo (posse,
indireto) --- esta última agora resultado de uma política de reposição
ótima, não um dado externo solto, em consistência com o restante do
formulário: a Parte IX, por exemplo, não aceita ``quantas instâncias de
recurso contratual existem'' como dado solto --- deriva isso de uma
condição de perpetuação aplicada sobre a demanda semanal já calculada. O
mesmo tratamento vale para consumíveis.

\hypertarget{x.1-aquisiuxe7uxe3o}{%
\subsubsection{X.1 Aquisição}\label{x.1-aquisiuxe7uxe3o}}

Separável por atividade sem ambiguidade. Por identidade direta
\(C_i^{\mathrm{V,op}} \equiv C_i^{\mathrm{V,cat}}\) (Parte VI):

\[\boxed{C^{\mathrm{cons}} = \sum_{i\in\mathcal{A}} \frac{Q_i\,\gamma_i^{\mathrm{V,cat}}}{\rho_i} = \sum_g c_g^{\mathrm{cons}}\, U_g}\qquad\text{[R\$ — direto, entra em $C^D$]}\]

\[C^{\mathrm{cons},\beta} = \sum_{i\in\mathcal{A}} \gamma_i^{\mathrm{V,cat},\beta}\cdot\frac{Q_i}{\rho_i}\]

Diferente do contratual, \(C^{\mathrm{cons}}\) não requer otimização de
cadeias --- a soma por atividade é exata, e coincide exatamente com a
produção realizada dado o fechamento de escopo da Parte IV. \(U_g\)
(unidades físicas do bem \(g\) usadas no projeto inteiro) é recuperado
com exatidão em §X.2, não apenas como total agregado.

\hypertarget{x.2-demanda-fuxedsica-derivada}{%
\subsubsection{X.2 Demanda física
derivada}\label{x.2-demanda-fuxedsica-derivada}}

O restante desta parte deriva \(U_g\) semana a semana --- necessário
para qualquer decisão de quando comprar, não apenas quanto. Para cada
atividade \(i\) e bem físico \(g\), a \textbf{taxa de consumo agregada}
do bem \(g\) por hora-equivalente de \(i\) soma os recursos consumíveis
de \(R_i\) que consomem \(g\) (§I.3.1) --- mesmo dado de cadastro que já
alimenta \(\gamma_i^{\mathrm{V,cat}}\) (Parte III), nenhum insumo novo:

\[\kappa_i^g \equiv \sum_{r\in R_i\,:\,g(r)=g} n_r\,\kappa_r \qquad [\text{unid}_g/\text{h}]\]

Usando a mesma razão \(q_{ij}/\rho_i\) que torna
\(C_{ij}^{\mathrm{V,op}}\) exata inclusive na semana de fechamento
(Parte V) --- ela já generaliza corretamente ``horas efetivas'' nas
semanas interiores e no residual ---, o consumo físico de \(g\) pela
atividade \(i\) na semana \(j\) segue a mesma construção, só trocando a
taxa monetária pela taxa física:

\[\boxed{D_{g,ij} = \kappa_i^g\cdot\frac{q_{ij}}{\rho_i}}\qquad[\text{unid}_g] \qquad\qquad \boxed{D_{g,j} = \sum_{i\,\in\,\mathcal{A}_j} D_{g,ij}}\qquad[\text{unid}_g]\]

com \(\mathcal{A}_j\) o conjunto de atividades ativas na semana \(j\)
(Parte IV). Nenhum cálculo novo além de agregação: \(q_{ij}\) e
\(\rho_i\) já são saída da Fase 3.

\textbf{Identidade de consistência.} Somando sobre toda a vida do
projeto e usando o fecho de escopo exato \(\sum_jq_{ij}=Q_i\):

\[\sum_j D_{g,j} = \sum_i\kappa_i^g\cdot\frac{Q_i}{\rho_i} \equiv U_g\]

recuperando exatamente o \(U_g\) de §X.1 --- a demanda física total já
era implicitamente conhecida pelo formulário original; faltava sua
distribuição semana a semana.

\hypertarget{x.3-balanuxe7o-de-estoque}{%
\subsubsection{X.3 Balanço de estoque}\label{x.3-balanuxe7o-de-estoque}}

Seja \(X_{g,j}\ge0\) a quantidade pedida do bem \(g\) na semana \(j\), e
\(L_g\ge0\) seu lead time de entrega (semanas, cadastral). A quantidade
recebida na semana \(j\) é \(O_{g,j}=X_{g,j-L_g}\), e a trajetória de
estoque satisfaz o balanço padrão:

\[\boxed{S_{g,j} = S_{g,j-1}+O_{g,j}-D_{g,j}}\qquad\text{s.a.}\quad S_{g,j}\ge0\ \ \forall j\]

com \(S_{g,j_{\min}-1}=S_g^0\) (tipicamente \(0\), salvo estoque de
mobilização declarado à parte). A restrição \(S_{g,j}\ge0\) é a condição
de \textbf{não-ruptura}: nenhuma semana consome mais do que já chegou.
Esta recursão conecta \(D_{g,j}\) (§X.2, já derivado) à decisão de
compra \(X_{g,j}\) (§X.4, a otimizar), produzindo \(S_{g,j}\) como saída
--- não mais como dado de entrada.

\hypertarget{x.4-poluxedtica-uxf3tima-de-reposiuxe7uxe3o-wagnerwhitin}{%
\subsubsection{X.4 Política ótima de reposição
(Wagner--Whitin)}\label{x.4-poluxedtica-uxf3tima-de-reposiuxe7uxe3o-wagnerwhitin}}

\(D_{g,j}\) é, por construção, não estacionária (perfis trapezoidais de
mobilização, atividades sobrepostas com datas distintas) e inteiramente
conhecida de antemão, já determinada na Fase 1 --- não é problema de
previsão sob incerteza, mas de lot-sizing determinístico, com solução
exata por programação dinâmica (Wagner \& Whitin, 1958). O problema
decompõe por bem \(g\), isoladamente --- mesma lógica de decomposição do
dimensionamento de pool (Parte IX/Fase 4), sem restrição de capital
compartilhado:

\[\min_{\{X_{g,j}\ge0\}}\ \sum_j\Big[c_g^{\text{ped}}\cdot\mathbb{1}[X_{g,j}>0]+c_g^{\text{hold}}\,S_{g,j}\Big]\qquad\text{s.a. balanço de §X.3},\ S_{g,j}\ge0\]

onde \(c_g^{\text{ped}}\) é o custo fixo de colocar um pedido (frete,
mínimo de fatura, administração de compra) --- cadastral por bem,
análogo em papel a \(c_r^{\text{mob}}+c_r^{\text{demob}}\) da Parte IX.

\textbf{Propriedade de estoque zero.} Existe solução ótima em que todo
pedido chega com estoque zero --- nunca é ótimo pedir enquanto ainda há
saldo positivo (resultado clássico de Wagner--Whitin). O horizonte de
demanda particiona-se em lotes consecutivos \([a,b]\): um único pedido,
colocado na semana \(a-L_g\), cobre a demanda de \(a\) até \(b\). O
estoque ao final da semana \(j\in[a,b]\) é o saldo ainda não consumido
dentro do lote, \(S_{g,j}=\sum_{k=j+1}^bD_{g,k}\); somando o custo de
posse ao longo do lote, cada \(D_{g,k}\) contribui \((k-a)\) vezes (uma
semana para cada \(j<k\) em que ainda estava em estoque):

\[\boxed{\mathrm{Custo}(a,b) = c_g^{\text{ped}} + c_g^{\text{hold}}\sum_{k=a+1}^{b}(k-a)\,D_{g,k}}\]

computável exatamente para qualquer \((a,b)\) candidato, sem depender de
lotes vizinhos --- mesma lógica de ``cada tentativa é recalculada do
zero, sem resíduo'' já usada na busca de \(N_i\) (§I.4).

\textbf{Recursão.} Seja \([j_{\min}^g,j_{\max}^g]\) o intervalo de
semanas com \(D_{g,j}>0\), e \(f(j_{\min}^g-1)\equiv0\):

\[\boxed{f(b) = \min_{a\,\in\,[j_{\min}^g,\,b]}\big[f(a-1)+\mathrm{Custo}(a,b)\big]}\]

O custo ótimo de reposição de \(g\) é \(f(j_{\max}^g)\); a partição
ótima em lotes --- logo as semanas de pedido e
\(X_{g,a-L_g}=\sum_{k=a}^bD_{g,k}\) --- é recuperada por backtracking do
argumento minimizador em cada passo. Complexidade \(O(n^2)\) no número
de semanas com demanda; refinamento \(O(n\log n)\) existe via teorema do
horizonte de planejamento (Wagelmans, Van Hoesel \& Kolen, 1992),
desnecessário na escala típica de um cronograma de obra. O lead time
\(L_g\) não altera a estrutura --- apenas desloca cada semana de pedido
\(a\) para \(a-L_g\); se isso cai antes do início do cronograma,
sinaliza necessidade de pedido de pré-mobilização (verificado em
§XIV.9).

Como o preço \(p_{g(r)}\) permanece constante (sem desconto por volume,
§X.6), o custo de \textbf{aquisição} \(C^{\mathrm{cons}}\) (§X.1) é
inalterado pela política de lotes --- depende só de \(U_g\), invariante
ao cronograma. A política de reposição determina exclusivamente a
repartição do custo de \textbf{posse} e introduz o custo de
\textbf{pedido}, nunca o total de material comprado.

\hypertarget{x.5-custo-de-posse-custo-de-pedido-e-atualizauxe7uxe3o-do-indireto}{%
\subsubsection{X.5 Custo de posse, custo de pedido e atualização do
indireto}\label{x.5-custo-de-posse-custo-de-pedido-e-atualizauxe7uxe3o-do-indireto}}

O custo de posse passa a ser saída da otimização, não soma sobre
\(S_{g,j}\) dado:

\[\boxed{C^{\mathrm{inv}} = \sum_g\Big(f(j_{\max}^g)-n_g^{\text{ped}}\cdot c_g^{\text{ped}}\Big)}\qquad\text{[R\$ — indireto, com $\beta^{\mathrm{IND}}$]}\]

com \(n_g^{\text{ped}}\) o número de lotes (pedidos) na partição ótima
de \(g\) --- subtrair seu custo de \(f(j_{\max}^g)\) isola exatamente a
parcela de posse. O custo de pedidos ganha categoria própria --- é
logística de suprimento, não posse, e carrega \(\beta^{\mathrm{IND}}\)
pela mesma razão que \(C^{MOB,eq}\) (Parte XII.5):

\[\boxed{C^{\mathrm{PED}} = \sum_g n_g^{\text{ped}}\cdot c_g^{\text{ped}}}\qquad\qquad C^{\mathrm{PED},\beta} = (1+\beta^{\mathrm{IND}})\cdot C^{\mathrm{PED}}\]

Os totais do indireto (Parte XII.6) passam a incluir esse termo:

\[C^I = \sum_{l\in\mathcal{L}} C_l^{I,o} + C^{MOB,eq,o} + C^{\mathrm{PED}} + C^{\mathrm{inv}}\]
\[C^{I,\beta} = \sum_{l\in\mathcal{L}} C_l^{I,\beta} + C^{MOB,eq,\beta} + C^{\mathrm{PED},\beta} + C^{\mathrm{inv}}\,(1+\beta^{\mathrm{IND}})\]

Nenhuma outra parte do formulário muda: \(C^D\), \(C^{\mathrm{cat}}\),
\(\mathrm{BDI}\), \(R\) mantêm suas fórmulas originais (Partes XI e
XIII), apenas com \(C^I\) recomputado.

\hypertarget{x.6-escopo}{%
\subsubsection{X.6 Escopo}\label{x.6-escopo}}

Ficam fora deste formulário, como extensões independentes: decomposição
de \(c_g^{\text{hold}}\) em capital/armazenagem/risco; desconto por
volume de compra (\(p_g\) permanece constante --- introduzi-lo trocaria
\(\mathrm{Custo}(a,b)\) por uma função não necessariamente separável,
exigindo variante do Wagner--Whitin com custo marginal decrescente);
estoque de segurança sob incerteza (a formulação é inteiramente
determinística, coerente com \(D_{g,j}\) já ser conhecida a partir do
cronograma resolvido --- introduzir variância exigiria propagar
incerteza de \(\eta_{p,j},\eta_{w,j}\) até \(D_{g,j}\), adendo de
natureza distinta); e perecibilidade/validade de lote. Cada um é
tratável como extensão adicional sobre a base agora estabelecida por
\(D_{g,j}\) e pela recursão de balanço (§X.2--X.3).

\begin{center}\rule{0.5\linewidth}{0.5pt}\end{center}

\hypertarget{parte-xi-custo-direto-total-do-projeto}{%
\subsection{Parte XI --- Custo Direto Total do
Projeto}\label{parte-xi-custo-direto-total-do-projeto}}

Esta seção apenas soma os dois grandes blocos de custo direto calculados
até aqui: o pool de recursos contratuais (Parte IX) e os consumíveis
(Parte X). Não há cálculo novo --- é a consolidação, em nível de
projeto, do que já foi apurado por atividade e por recurso.

\[C^D = C^{\mathrm{pool}}+C^{\mathrm{cons}}\qquad\text{[R\$ — sem BDI]}\]
\[C^{D,\beta} = C^{\mathrm{pool},\beta}+C^{\mathrm{cons},\beta}\qquad\text{[R\$ — com BDI; nível projeto, distinto de }C_i^{D,\beta}\text{ (atividade)}\]
\[\beta^D = C^{D,\beta}/C^D-1\]

\(\beta^D\) é o BDI médio implícito nesse bloco direto --- um
diagnóstico intermediário útil antes de somar os custos indiretos na
Parte XIII.

\begin{center}\rule{0.5\linewidth}{0.5pt}\end{center}

\hypertarget{parte-xii-custos-indiretos}{%
\subsection{Parte XII --- Custos
Indiretos}\label{parte-xii-custos-indiretos}}

Custos indiretos são aqueles que existem para viabilizar o projeto como
um todo, não para produzir uma atividade específica --- administração
central, topografia, laboratório, QSMS, mobilização de canteiro.
Diferente do custo direto, que só existe enquanto uma atividade está
efetivamente em produção, a maior parte do indireto é cobrada pelo
simples fato de o projeto estar em andamento, semana a semana, do início
ao fim do cronograma.

\hypertarget{xii.1-estrutura-categorias-e-drivers}{%
\subsubsection{XII.1 Estrutura: categorias e
drivers}\label{xii.1-estrutura-categorias-e-drivers}}

\[\mathcal{L} = \{ADM, TOP, LAB, QSMS, MOB^{\text{site}}, \ldots\}\quad\text{(conjunto extensível)}\]

\[\text{driver}(l)\in\{\texttt{perm},\ \texttt{onetime}\}\]

\begin{itemize}
\tightlist
\item
  \texttt{perm}: custo mensal fixo enquanto o projeto está mobilizado
  (ADM, TOP, LAB, QSMS)
\item
  \texttt{onetime}: evento único de borda --- mobilização/desmobilização
  de canteiro (\(MOB^{\text{site}}\))
\end{itemize}

Em outras palavras: categorias \texttt{perm} (como administração
central) geram custo toda semana em que o projeto está ativo,
proporcional ao makespan do cronograma; categorias \texttt{onetime}
(como mobilização de canteiro) geram custo pontual, uma vez no início e
uma vez no fim, independentemente de quanto tempo o projeto durar.

\(n = j^{\max}-j^{\min}\) (makespan do projeto em semanas).
\(j^{\min}=\min_i j_i^0\), \(j^{\max}=\max_i(j_i^0+N_i)\).

\hypertarget{xii.2-entrada-por-categoria}{%
\subsubsection{XII.2 Entrada por
categoria}\label{xii.2-entrada-por-categoria}}

\[l = \Bigl(\text{driver}(l),\ \{\Gamma_l^{k,\text{mensal}}\}_{k\in\{MO,EQ,COMB\}},\ (c_l^{\text{mob}},c_l^{\text{demob}})\ \text{ou}\ \Gamma_l^{VM,\text{mensal}}\Bigr)\]

\hypertarget{xii.3-custo-semanal-por-sub-categoria}{%
\subsubsection{XII.3 Custo semanal por
sub-categoria}\label{xii.3-custo-semanal-por-sub-categoria}}

\textbf{Driver \texttt{perm}:}

\[\gamma_l^{k,\text{sem}} = \frac{\Gamma_l^{k,\text{mensal}}}{s_{\text{mes}}}, \qquad m_l(j)=1 \ \ \forall j\in\{0,\ldots,n-1\}\]

\textbf{Driver \texttt{onetime}:}

\[\gamma_l^{VM,\text{sem}}(j) = c_l^{\text{mob}}\,\mathbb{1}[j=j^{\min}] + c_l^{\text{demob}}\,\mathbb{1}[j=j^{\max}-1]\]

\emph{(para \(n=1\): mob e demob caem na mesma semana, somando-se)}

\[C_{lj}^{I,o} = \Bigl(\gamma_l^{MO,\text{sem}}+\gamma_l^{EQ,\text{sem}}+\gamma_l^{COMB,\text{sem}}+\gamma_l^{VM,\text{sem}}(j)\Bigr)\cdot m_l(j)\]

\hypertarget{xii.4-camada-bdi-do-indireto}{%
\subsubsection{XII.4 Camada BDI do
indireto}\label{xii.4-camada-bdi-do-indireto}}

\[\boxed{C_{lj}^{I,\beta} = \Bigl[\gamma_l^{MO,\text{sem}}(1+\beta^{MO})+\gamma_l^{EQ,\text{sem}}(1+\beta^{EQ})+\gamma_l^{COMB,\text{sem}}(1+\beta^{COMB})+\gamma_l^{VM,\text{sem}}(j)(1+\beta^{IND})\Bigr]\cdot m_l(j)}\]

\textbf{Totais por categoria:}

\[C_l^{I,o} = \sum_j C_{lj}^{I,o}\]
\[C_l^{I,\beta} = \Gamma_l^{MO}(1+\beta^{MO})+\Gamma_l^{EQ}(1+\beta^{EQ})+\Gamma_l^{COMB}(1+\beta^{COMB})+\Gamma_l^{VM}(1+\beta^{IND})\]

\hypertarget{xii.5-mobilizauxe7uxe3o-do-pool-de-equipamento-indireta-por-cadeia}{%
\subsubsection{XII.5 Mobilização do pool de equipamento (indireta, por
cadeia)}\label{xii.5-mobilizauxe7uxe3o-do-pool-de-equipamento-indireta-por-cadeia}}

\[\boxed{C^{MOB,eq,o} = \sum_r N_{\text{chains}}^r\cdot(c_r^{\text{mob}}+c_r^{\text{demob}})}, \qquad C^{MOB,eq,\beta} = (1+\beta^{IND})\cdot C^{MOB,eq,o}\]

\hypertarget{xii.6-totais-do-indireto}{%
\subsubsection{XII.6 Totais do
indireto}\label{xii.6-totais-do-indireto}}

\[C^I = \sum_{l\in\mathcal{L}} C_l^{I,o} + C^{MOB,eq,o} + C^{\mathrm{inv}}\]
\[C^{I,\beta} = \sum_{l\in\mathcal{L}} C_l^{I,\beta} + C^{MOB,eq,\beta} + C^{\mathrm{inv}}\,(1+\beta^{\mathrm{IND}})\]
\[\bar\beta^I = C^{I,\beta}/C^I-1\]

\begin{center}\rule{0.5\linewidth}{0.5pt}\end{center}

\hypertarget{parte-xiii-suxedntese-preuxe7o-total-e-viabilidade}{%
\subsection{Parte XIII --- Síntese: Preço Total e
Viabilidade}\label{parte-xiii-suxedntese-preuxe7o-total-e-viabilidade}}

Esta seção final reúne tudo que foi calculado: soma-se o custo direto e
indireto, com e sem BDI, compara-se o resultado com o envelope de
catálogo, e verifica-se se a proposta comercial resultante é viável
frente ao piso regulatório.

\textbf{Custo de catálogo} (agregado por atividade):

\[C^{\mathrm{cat}} = \sum_{i\in\mathcal{A}} C_i^{\mathrm{cat}} = \sum_{i\in\mathcal{A}} \frac{\gamma_i^{\mathrm{cat}}\,Q_i}{\rho_i}\qquad\text{[R\$ — envelope de catálogo, base de precificação]}\]

\[C^{\mathrm{tot}} = C^D+C^I \qquad\text{[R\$ — custo operacional total, sem BDI]}\]

\textbf{Preço de venda e BDI implícito.} O preço é a soma dos blocos
direto e indireto já com a camada granular de \(\beta\) aplicada (Partes
VIII e XII):

\[P \equiv C^{D,\beta}+C^{I,\beta} \qquad\text{[R\$ — preço de venda]}\]

O BDI não é um segundo parâmetro de markup somado independentemente
sobre \(C^{\mathrm{cat}}\): é a taxa agregada, \textbf{derivada},
resultante de aplicar a camada granular de \(\beta\)'s --- uma por
categoria de recurso
(\(\beta^{EQ},\beta^{MO},\beta^{MAT},\beta^{COMB},\beta^{IND}\)) --- e
agregá-la ao nível do envelope de catálogo:

\[\boxed{\mathrm{BDI} \equiv \frac{P}{C^{\mathrm{cat}}}-1}\]

Não há dupla contagem: existe uma única camada de markup, aplicada item
a item nas Partes VIII e XII; o BDI é apenas a leitura agregada dessa
mesma camada, útil para comparação com o piso regulatório. Categorias de
custo indireto que fisicamente compõem a estrutura clássica de BDI
(administração central, canteiro) já estão computadas dentro de
\(C^{I,\beta}\) com seu próprio \(\beta^{IND}\) --- não devem ser
adicionadas novamente como percentual do BDI agregado;
\(\mathrm{BDI}_{\mathrm{TCU,min}}\) funciona como \textbf{piso de
verificação} sobre a taxa agregada resultante, não como um segundo
insumo a ser somado.

\[R = C^{\mathrm{tot}}/C^{\mathrm{cat}}, \qquad \beta^{\mathrm{eq}} = P/C^{\mathrm{tot}}-1, \qquad \mu = (P-C^{\mathrm{tot}})/P\]

\textbf{Condição de viabilidade:}

\[\boxed{\mathrm{BDI} \;\ge\; \mathrm{BDI}_{\mathrm{TCU,min}}}\]

\(\mathrm{BDI}_{\mathrm{TCU,min}}\) é determinado pela fórmula
multiplicativa TCU (overhead não capturado nos \(\beta\)'s granulares,
lucro mínimo e tributos) --- é o único piso de verificação que este
formulário exige.

Poderia parecer que uma segunda verificação também é necessária:
garantir que o preço pelo menos cubra o custo realmente incorrido, isto
é, \(\mathrm{BDI}\ge R-1\) (equivalente a \(P\ge C^{\mathrm{tot}}\)).
Essa condição, porém, não é uma segunda verificação independente --- é
consequência aritmética automática de \(\mathrm{BDI}\) e \(R\), ambos já
calculados, não um novo insumo a se checar à parte. Substituindo
\(P=C^{\mathrm{cat}}(1+\mathrm{BDI})\) e
\(C^{\mathrm{tot}}=R\cdot C^{\mathrm{cat}}\) em
\(\beta^{\mathrm{eq}}=P/C^{\mathrm{tot}}-1\):

\[\beta^{\mathrm{eq}} = \frac{C^{\mathrm{cat}}(1+\mathrm{BDI})}{R\cdot C^{\mathrm{cat}}}-1 = \frac{1+\mathrm{BDI}}{R}-1\]

logo
\(\beta^{\mathrm{eq}}>0 \iff 1+\mathrm{BDI}>R \iff \mathrm{BDI}>R-1\)
--- as duas expressões descrevem exatamente a mesma coisa, apenas
escritas de formas diferentes. Verificar
\(\mathrm{BDI}\ge\mathrm{BDI}_{\mathrm{TCU,min}}\) já é suficiente; não
há necessidade de comparar \(\mathrm{BDI}\) com \(R-1\) à parte. \(R\)
permanece útil como \textbf{diagnóstico}, não como insumo de decisão:
evidencia quanto do custo provém de ineficiência prevista, ajudando a
explicar por que o BDI resultante ficou alto ou baixo --- não substitui
nem complementa a verificação regulatória.

\begin{center}\rule{0.5\linewidth}{0.5pt}\end{center}

\hypertarget{parte-xiv-algoritmo}{%
\subsection{Parte XIV --- Algoritmo}\label{parte-xiv-algoritmo}}

O algoritmo a seguir é a receita passo a passo que implementa todas as
fórmulas das partes anteriores, na ordem correta de dependência: cada
fase usa apenas resultados já calculados pelas fases anteriores, nunca
antecipa um valor que ainda não existe. Quem já entendeu as Partes I a
XIII pode ler esta seção como um mapa de onde cada fórmula entra no
processo.

O algoritmo opera em \textbf{10 fases sequenciais} (incluindo a Fase 4b,
dinâmica de estoque e reposição). As Fases 0 e 1 requerem solvers
externos; as demais são computações exatas dadas as saídas anteriores
--- a própria Fase 4b resolve um problema de programação dinâmica de
solução exata conhecida (Wagner--Whitin), não uma heurística.

\begin{verbatim}
══════════════════════════════════════════════════════════════════════
ENTRADA
══════════════════════════════════════════════════════════════════════

  A = {A_i}_{i=1}^{n_ativ}        — conjunto de atividades

  Para cada i:
    Q_i                            — quantitativo total [un]
    G_i = (V_i, A_i^syn, cap)      — grafo de sinergia da equipe
    R_i = {r: (cat_r, n_r, V_r^mensal ou (κ_r, p_{g(r)}))}
    W_i = (m_i^0, μ_i^+, m_i^max, μ_i^-, m_i^1)
    B_i, θ_i: B_i → (0,1]
    (h_i, d_i)                     — regime de jornada

  Parâmetros globais:
    β^EQ, β^MO, β^MAT, β^COMB, β^IND, s_mes
    {η_{p,j}, η_{w,j}}              — calendário global de desempenho e clima
    // η_{s,ij} é derivado por atividade a partir de G_i e do perfil {m_ij} (Fase 1)
    // ε_{ij} = η_{s,ij} · η_{p,j} · η_{w,j}

  Para cada recurso r:
    c_r^mob, c_r^demob             — custos de mobilização/desmobilização
    G_r^viável                     — gap máximo absoluto [semanas]
    // c_r^standby ≡ c_r^sem: pool é custo de permanência puro (§IX.3)

  Categorias indiretas:
    {(l, driver(l), {Γ_l^{k,mensal}}, c_l^mob, c_l^demob)}_{l ∈ L}

  Para cada bem físico g (consumível, cf. g(r) em R_i):
    c_g^hold, c_g^ped, L_g          — custo de posse [R$/unid_g/sem], custo de pedido [R$], lead time [sem]
    j^{procure,min}                 — semana mais cedo em que um pedido pode ser colocado (opcional)
\end{verbatim}

\hypertarget{xiv.0-fase-0-resoluuxe7uxe3o-de-produtividade}{%
\subsubsection{XIV.0 --- Fase 0: Resolução de
Produtividade}\label{xiv.0-fase-0-resoluuxe7uxe3o-de-produtividade}}

Em termos simples: aqui se calcula quanto cada equipe consegue produzir
por hora, a partir do grafo de sinergia ou de um valor calibrado
externamente.

\begin{verbatim}
Para cada i ∈ A:
  Se i.grafo_sinergia está presente:
    Identificar topologia de G_i:
      série    → ρ_i = min_{(u,v) ∈ caminho_crítico} cap(u,v)
      paralelo → ρ_i = Σ_{(s_i,v) ∈ A_i^syn} cap(s_i, v)
      misto    → ρ_i = max-fluxo(s_i → t_i, G_i^D)  [Ford-Fulkerson]
    Armazenar ρ_i como producao_horaria  [cache auditável]
  Senão:
    ρ_i ← i.producao_horaria             [entrada direta calibrada externamente]

  Validar ρ_i > 0  [equipe inviável se algum corte tem capacidade zero]

SAÍDA: {ρ_i} ∀i
\end{verbatim}

\hypertarget{xiv.1-fase-1-cronograma-interface-com-solver-rcpsp-externo}{%
\subsubsection{XIV.1 --- Fase 1: Cronograma (interface com solver RCPSP
externo)}\label{xiv.1-fase-1-cronograma-interface-com-solver-rcpsp-externo}}

Em termos simples: aqui se determina quantas semanas cada atividade vai
durar e como as frentes de trabalho se distribuem ao longo do tempo, até
fechar exatamente o quantitativo contratado.

\begin{verbatim}
Pré-condição: {ρ_i} disponível ∀i

Para cada i em ordem topológica de B_i:

  // j_i^0 depende só de predecessores, já resolvidos — calculado antes da busca por N_i,
  // pois a busca abaixo precisa de semanas absolutas j para consultar η_{p,j}, η_{w,j}
  j_i^0 ← max_{b ∈ B_i} min{j : Φ_b(j) ≥ θ_i(b)}  [restrição exata de quantidade]
    com Φ_b(j) = (1/Q_b) · Σ_{j'=j_b^0}^{j} q_{b,j'}
    // q_{b,j'} vem da DistribuicaoSemanal já calculada dos predecessores b, já com
    // fechamento fracionário aplicado

  N_i ← max(1, ⌈Q_i / (ρ_i · h_i · d_i · m_i^max)⌉)   // chute inicial (eficiência plena)

  Repetir:
    {m_i(τ)}_{τ=0}^{N_i-1} ← RampProfile.build(N_i, m_i^0, μ_i^+, m_i^max, μ_i^-, m_i^1)
    // perfil inteiramente determinado pelo N_i testado — nenhuma incógnita nesta avaliação

    Para τ = 0, ..., N_i-1:
      j ← j_i^0 + τ                                    // semana absoluta, para consultar o calendário
      Se i.grafo_sinergia presente:
        η_{s,iτ} ← max-fluxo(G_i^D | m_i(τ)) / max-fluxo(G_i^D | m_i^max)
      Senão:
        η_{s,iτ} ← m_i(τ) / m_i^max
      ε_{iτ} ← η_{s,iτ} · η_{p,j} · η_{w,j}

    prod ← Σ_τ ρ_i · ε_{iτ} · h_i · d_i · m_i(τ)

    Se prod ≥ Q_i:
      convergiu; {m_{ij}} ← {m_i(τ)}; {ε_{ij}} ← {ε_{iτ}}
    Senão:
      Se ε_{iτ} = 0 para toda semana τ futura alcançável dentro do calendário disponível:
        Marcar A_i INVIÁVEL; interromper busca (→ Fase 9, NOGO)
      Senão:
        N_i ← N_i + 1
        // Π(N) = prod acumulado é estritamente crescente em N (ver §I.4): a cada
        // incremento, as semanas já existentes só podem crescer (o piso de ramp-down
        // relaxa) e a semana nova soma ≥ ρ_i · ε_{i,N-1} · h_i · d_i · m_i^1 > 0 —
        // convergência garantida sempre que o calendário não anula ε indefinidamente.
        // Cada tentativa é recalculada do zero a partir da fórmula de m_i(τ;N); nenhum
        // valor de uma tentativa rejeitada é reaproveitado na seguinte.

  // Fechamento exato de escopo — última semana ativa (τ relativo, 0-indexado)
  τ^1 ← N_i − 1
  q_residual ← Q_i − Σ_{τ=0}^{τ^1-1} ρ_i · ε_{i,τ} · h_i · d_i · m_i(τ)
  φ_i ← q_residual / (ρ_i · ε_{i,τ^1} · h_i · d_i · m_i(τ^1))

  Verificar: r + d ≤ N_i  [consistência com W_i]

Armazenar (j_i^0, N_i, {m_{ij}}, {ε_{ij}}, φ_i)

SAÍDA: {j_i^0, N_i, {m_{ij}}, {ε_{ij}}, φ_i} ∀i
\end{verbatim}

\hypertarget{xiv.2-fase-2-taxas-horuxe1rias-de-catuxe1logo}{%
\subsubsection{XIV.2 --- Fase 2: Taxas Horárias de
Catálogo}\label{xiv.2-fase-2-taxas-horuxe1rias-de-catuxe1logo}}

Em termos simples: aqui se traduz o cadastro de cada recurso --- aluguel
mensal, taxa de consumo, preço de insumo --- em um custo por hora, e
depois esse custo é agregado por equipe.

\begin{verbatim}
Para cada i ∈ A:
  Para cada r ∈ R_i:
    Se cat(r) ∈ {EQ, MO}:
      c_r ← V_r^mensal / (s_mes · d_i · h_i)          [R$/h]
    Se cat(r) ∈ {MAT, COMB}:
      c_r ← κ_r · p_{g(r)}                             [R$/h_op, agnóstico em unid_g]

  Para k ∈ {EQ, MO, MAT, COMB}:
    γ_i^k ← Σ_{r: cat(r)=k} c_r · n_r

  γ_i^{K,cat} ← γ_i^EQ + γ_i^MO
  γ_i^{V,cat} ← γ_i^MAT + γ_i^COMB
  γ_i^{cat}   ← γ_i^{K,cat} + γ_i^{V,cat}
  f_i^V       ← γ_i^{V,cat} / γ_i^{cat}
  p_i^{cat}   ← γ_i^{cat} / ρ_i                       [R$/un — PU de catálogo]

SAÍDA: {γ_i^k, γ_i^{K,cat}, γ_i^{V,cat}, γ_i^{cat}, f_i^V, p_i^{cat}} ∀i
\end{verbatim}

\hypertarget{xiv.3-fase-3-quantitativos-e-custos-semanais-diretos}{%
\subsubsection{XIV.3 --- Fase 3: Quantitativos e Custos Semanais
Diretos}\label{xiv.3-fase-3-quantitativos-e-custos-semanais-diretos}}

Em termos simples: aqui se calcula, semana a semana, quanto cada
atividade produziu e quanto isso custou --- e se confere que a soma bate
exatamente com o quantitativo contratado.

\begin{verbatim}
Para cada i ∈ A, para cada j ∈ [j_i^0, j_i^0 + N_i - 1]:
  Se j < j_i^0 + N_i - 1:
    q_{ij} ← ρ_i · ε_{ij} · h_i · d_i · m_{ij}          [semana interior]
  Senão:
    q_{ij} ← q_residual                                  [semana de fechamento, da Fase 1]

  C_{ij}^{K,op}   ← γ_i^{K,cat} · h_i · d_i · m_{ij}       [nominal — não prorateado]
  C_{ij}^{V,op}   ← γ_i^{V,cat} · q_{ij} / ρ_i             [prorateado exatamente pelo quantitativo]
  C_{ij}^{op}     ← C_{ij}^{K,op} + C_{ij}^{V,op}

Para cada i:
  M_i          ← Σ_j m_{ij}
  C_i^{K,op}  ← γ_i^{K,cat} · h_i · d_i · M_i
  C_i^{V,op}  ← γ_i^{V,cat} · Q_i / ρ_i           [= Σ_j C_{ij}^{V,op}, exato]
  C_i^{op}    ← C_i^{K,op} + C_i^{V,op}
  p_i^{op}    ← C_i^{op} / Q_i                      [R$/un]

Verificação obrigatória:
  ∀i: Σ_j q_{ij} = Q_i                               [fecho de escopo exato]

C_i^{cat} ← γ_i^{cat} · Q_i / ρ_i
C^{cat}   ← Σ_i C_i^{cat}                            // envelope de catálogo

SAÍDA: {q_{ij}, C_{ij}^{op}, C_i^{op}, M_i, p_i^{op}, C_i^{cat}} ∀i
       C^{cat}
\end{verbatim}

\hypertarget{xiv.4-fase-4-dimensionamento-e-custo-do-pool}{%
\subsubsection{XIV.4 --- Fase 4: Dimensionamento e Custo do
Pool}\label{xiv.4-fase-4-dimensionamento-e-custo-do-pool}}

Em termos simples: aqui se decide, para cada tipo de recurso contratual,
quantas instâncias contratar e quando vale a pena mantê-las paradas
entre duas atividades em vez de desmobilizar e remobilizar.

\begin{verbatim}
// Decomposição por nível (histograma): para um tipo fungível homogêneo, a alocação
// ótima decompõe-se em decisões independentes por "fatia" de altura unitária — cada
// fatia corresponde a uma cadeia de perpetuação da Parte IX. Ref.: SDVSP — Freling,
// Wagelmans & Paixão (2001).

C^{pool} ← 0; C^{pool,β} ← 0; C^{MOB,eq} ← 0; C^{MOB,eq,β} ← 0

Para cada tipo τ ∈ T:
  n_j(τ) ← Σ_{i: j ∈ [j_i^0, j_i^0+N_i-1]} m_{ij} · n_τ^i   ∀j
  n_max(τ) ← max_j n_j(τ)
  G_τ^* ← (c_τ^mob + c_τ^demob) / c_τ^sem

  sem_τ ← 0; trips_τ ← 0

  Para l = 1, 2, ..., n_max(τ):                       // uma "fatia" por unidade de fleet
    on_l ← {j : n_j(τ) ≥ l}                             // semanas em que a l-ésima instância é exigida
    blocos ← intervalos maximais consecutivos de on_l
    Fundir blocos consecutivos [b_k, b_{k+1}] cujo gap = início(b_{k+1}) − fim(b_k)
      satisfaça gap ≤ min(G_τ^*, G_τ^viável)            // perpetuação da instância desta fatia
    Para cada bloco fundido (após a fusão):
      sem_τ   ← sem_τ + duração(bloco)                  // inclui gaps internos perpetuados
      trips_τ ← trips_τ + 1                               // um par mob/demob por bloco fundido

  C^{pool}     ← C^{pool}     + sem_τ · c_τ^sem
  C^{pool,β}   ← C^{pool,β}   + (1 + β_{cat(τ)}) · sem_τ · c_τ^sem
  C^{MOB,eq}   ← C^{MOB,eq}   + trips_τ · (c_τ^mob + c_τ^demob)
  C^{MOB,eq,β} ← C^{MOB,eq,β} + (1 + β^IND) · trips_τ · (c_τ^mob + c_τ^demob)

SAÍDA: C^{pool}, C^{pool,β}, C^{MOB,eq}, C^{MOB,eq,β}
\end{verbatim}

\hypertarget{xiv.4b-fase-4b-dinuxe2mica-de-estoque-e-reposiuxe7uxe3o-wagnerwhitin}{%
\subsubsection{XIV.4b --- Fase 4b: Dinâmica de Estoque e Reposição
(Wagner--Whitin)}\label{xiv.4b-fase-4b-dinuxe2mica-de-estoque-e-reposiuxe7uxe3o-wagnerwhitin}}

Em termos simples: aqui se decide, para cada bem consumível, quando e
quanto comprar, cobrindo exatamente a demanda derivada do cronograma ao
menor custo combinado de pedido e posse.

\begin{verbatim}
Pré-condição: {q_ij, ρ_i} disponíveis (Fase 3); {κ_i^g} do cadastro (Fase 2, via R_i)

Para cada bem físico g:

  // Demanda derivada — nenhum novo insumo, só agregação (§X.2)
  Para cada j:
    D_{g,j} ← Σ_{i ∈ A_j} κ_i^g · q_{ij} / ρ_i

  [j_min^g, j_max^g] ← intervalo de semanas com D_{g,j} > 0

  // Programação dinâmica (Wagner-Whitin) — solução exata, não heurística
  f(j_min^g - 1) ← 0
  Para b = j_min^g, ..., j_max^g:
    f(b) ← min_{a ∈ [j_min^g, b]} [ f(a-1) + c_g^ped + c_g^hold · Σ_{k=a+1}^{b} (k-a)·D_{g,k} ]
    a*(b) ← argmin correspondente

  // Backtracking — recupera a partição ótima em lotes
  lotes_g ← {}; b ← j_max^g
  Enquanto b ≥ j_min^g:
    a ← a*(b)
    lotes_g ← lotes_g ∪ {[a, b]}
    X_{g, a - L_g} ← Σ_{k=a}^{b} D_{g,k}          // pedido colocado L_g semanas antes de a
    b ← a - 1

  // Reconstrução da trajetória de estoque (§X.4)
  Para cada lote [a,b] ∈ lotes_g, para j = a, ..., b:
    S_{g,j} ← Σ_{k=j+1}^{b} D_{g,k}

  n_g^ped ← |lotes_g|
  C_g^{inv} ← c_g^hold · Σ_j S_{g,j}
  C_g^{PED} ← n_g^ped · c_g^ped

  Verificar: ∀ lote [a,b] ∈ lotes_g: a - L_g ≥ j_min   [senão, sinalizar pedido de pré-mobilização]

C^{inv} ← Σ_g C_g^{inv}
C^{PED} ← Σ_g C_g^{PED}
C^{PED,β} ← (1 + β^IND) · C^{PED}

SAÍDA: {D_{g,j}, S_{g,j}, X_{g,j}} ∀g,j ; C^{inv}, C^{PED}, C^{PED,β}
\end{verbatim}

\hypertarget{xiv.5-fase-5-estoque-consumuxedveis}{%
\subsubsection{XIV.5 --- Fase 5: Estoque
(Consumíveis)}\label{xiv.5-fase-5-estoque-consumuxedveis}}

Em termos simples: aqui se soma o custo de material e combustível
efetivamente consumido por todas as atividades do projeto.

\begin{verbatim}
C^{cons} ← Σ_i γ_i^{V,cat} · Q_i / ρ_i
  // Exato por construção: Σ_i γ_i^{V,cat}·Q_i/ρ_i ≡ Σ_i C_i^{V,op} (fecho de escopo, Parte IV)

Para cada i:
  γ_i^{V,cat,β} ← (1+β^MAT) · γ_i^MAT + (1+β^COMB) · γ_i^COMB

C^{cons,β} ← Σ_i γ_i^{V,cat,β} · Q_i / ρ_i

// C^{inv} e C^{PED} já disponíveis da Fase 4b — não recalculados aqui

SAÍDA: C^{cons}, C^{cons,β}
\end{verbatim}

\hypertarget{xiv.6-fase-6-custo-direto-total}{%
\subsubsection{XIV.6 --- Fase 6: Custo Direto
Total}\label{xiv.6-fase-6-custo-direto-total}}

Em termos simples: aqui se somam os dois blocos calculados nas fases
anteriores --- pool de recursos e consumíveis --- para chegar ao custo
direto do projeto inteiro.

\begin{verbatim}
C^D    ← C^{pool} + C^{cons}
C^{D,β} ← C^{pool,β} + C^{cons,β}
β^D    ← C^{D,β} / C^D − 1

SAÍDA: C^D, C^{D,β}, β^D
\end{verbatim}

\hypertarget{xiv.7-fase-7-custos-indiretos}{%
\subsubsection{XIV.7 --- Fase 7: Custos
Indiretos}\label{xiv.7-fase-7-custos-indiretos}}

Em termos simples: aqui se calculam os custos que existem para
viabilizar o projeto como um todo --- administração, canteiro,
mobilização de equipamento --- a partir do makespan do cronograma.

\begin{verbatim}
j_min ← min_i j_i^0
j_max ← max_i (j_i^0 + N_i)
n     ← j_max − j_min    [makespan em semanas]

Para cada l ∈ L:
  Se driver(l) = 'perm':
    Para k ∈ {MO, EQ, COMB}:
      Γ_l^k ← Γ_l^{k,mensal} · (n / s_mes)
    Γ_l^VM ← 0

  Se driver(l) = 'onetime':
    Γ_l^VM ← c_l^mob + c_l^demob
    Γ_l^k  ← 0  para k ∈ {MO, EQ, COMB}

  C_l^{I,o} ← Σ_k Γ_l^k + Γ_l^VM
  C_l^{I,β} ← Γ_l^MO·(1+β^MO) + Γ_l^EQ·(1+β^EQ)
               + Γ_l^COMB·(1+β^COMB) + Γ_l^VM·(1+β^IND)

// C^{inv} e C^{PED} já calculados na Fase 4b — apenas consumidos aqui
C^I    ← Σ_l C_l^{I,o} + C^{MOB,eq,o} + C^{PED} + C^{inv}
C^{I,β} ← Σ_l C_l^{I,β} + C^{MOB,eq,β} + C^{PED,β} + C^{inv}·(1+β^IND)
β^I    ← C^{I,β} / C^I − 1

SAÍDA: C^I, C^{I,β}, β^I
\end{verbatim}

\hypertarget{xiv.8-fase-8-suxedntese-e-preuxe7o-total}{%
\subsubsection{XIV.8 --- Fase 8: Síntese e Preço
Total}\label{xiv.8-fase-8-suxedntese-e-preuxe7o-total}}

Em termos simples: aqui se reúnem todos os blocos de custo em um preço
de venda final, e se calcula o BDI agregado resultante.

\begin{verbatim}
C^{tot} ← C^D + C^I
P       ← C^{D,β} + C^{I,β}
BDI     ← P / C^{cat} − 1                          // taxa agregada derivada, não insumo livre
R       ← C^{tot} / C^{cat}                          // razão custo real vs catálogo
β^eq   ← P / C^{tot} − 1
μ       ← (P − C^{tot}) / P

SAÍDA:
  C^D, C^I, C^{tot}       — custos sem BDI (base operacional)
  C^{D,β}, C^{I,β}, P    — preços com BDI (base de venda)
  BDI, R, β^eq, μ         — indicadores  [verificar BDI ≥ BDI_TCU_min; R = diagnóstico]
\end{verbatim}

\hypertarget{xiv.9-fase-9-validauxe7uxe3o-e-diagnuxf3stico}{%
\subsubsection{XIV.9 --- Fase 9: Validação e
Diagnóstico}\label{xiv.9-fase-9-validauxe7uxe3o-e-diagnuxf3stico}}

Em termos simples: aqui se checa se todos os resultados fazem sentido
--- se o escopo fechou exatamente, se o BDI atinge o piso regulatório,
se nenhum recurso ficou parado o tempo todo --- antes de aprovar a
proposta.

\begin{verbatim}
Verificações obrigatórias:
  [V1] ∀i: Σ_j q_{ij} = Q_i                           — fecho de escopo exato
  [V2] BDI ≥ BDI_TCU_min                                — cobertura mínima de overhead+lucro+tributos
  [V3] ∀r: f_standby(r) < 1                             — nenhum recurso 100% standby
  [V4] P > 0  e  C^{tot} > 0                            — consistência numérica
  [V5] ∀i: ρ_i > 0                                     — nenhuma equipe com gargalo de capacidade zero
  [V6] ∀i: busca de N_i convergiu sem sinalização de inviabilidade (Fase 1) — calendário viável
  [V7] ∀g,j: S_{g,j} ≥ 0                                — nenhuma ruptura de estoque na solução ótima (§X.3)
  [V8] ∀g: Σ_j D_{g,j} = U_g                            — consistência entre demanda física e consumo total (§X.2)
  [V9] ∀g, ∀ lote [a,b] ∈ lotes_g: a − L_g ≥ j_min      — viabilidade de prazo de pedido (senão, pré-mobilização)

Diagnósticos opcionais:
  [D1] taxa_util_τ = Σ_j n_j(τ) / (n · n_max(τ)) por tipo τ
         → próximo de 1: pool bem dimensionado; baixo: excesso de standby
  [D2] R_i = C_i^{op} / C_i^{cat} = f_i^V + (1 − f_i^V) · (1/ε̄_i)   [Parte VI/VIII]
         → razão custo operacional / catálogo por atividade
  [D2'] Δ_i^{cat} = C_i^{op} − C_i^{cat} = (R_i−1)·C_i^{cat}   [R$, Parte VI]
         → desvio de catálogo por atividade, em valor monetário
  [D3] {p_i^{op}} por atividade — custo unitário operacional [R$/un]
  [D4] Proporção direto/indireto: C^D/C^{tot}
  [D5] β_i^{eq} = (1−δ_i)·β_i^K + δ_i·β_i^V   — BDI equivalente por equipe
  [D6] n_g^ped e tamanho médio de lote por bem — fragmentação da política de reposição:
         muitos pedidos pequenos (custo de pedido domina) vs. poucos grandes (posse domina)

Saída final:
  GO   se [V1]–[V9] aprovadas
  NOGO + lista de violações  caso contrário
\end{verbatim}

\begin{center}\rule{0.5\linewidth}{0.5pt}\end{center}

\hypertarget{convenuxe7uxf5es-de-notauxe7uxe3o-derivado}{%
\subsection{§0 --- Convenções de Notação
{[}DERIVADO{]}}\label{convenuxe7uxf5es-de-notauxe7uxe3o-derivado}}

Antes de qualquer fórmula, vale entender como os símbolos são montados:
cada grandeza carrega até três informações no seu sobrescrito, o que
evita ambiguidade sem exigir nomes de variáveis longos. Uma vez
internalizado esse código, o resto do formulário se lê quase como texto
corrido.

\textbf{Regra de três slots} no superscrito. Nenhum símbolo carrega dois
papéis.

\begin{longtable}[]{@{}
  >{\raggedright\arraybackslash}p{(\columnwidth - 4\tabcolsep) * \real{0.3333}}
  >{\raggedright\arraybackslash}p{(\columnwidth - 4\tabcolsep) * \real{0.3333}}
  >{\raggedright\arraybackslash}p{(\columnwidth - 4\tabcolsep) * \real{0.3333}}@{}}
\toprule\noalign{}
\begin{minipage}[b]{\linewidth}\raggedright
Slot
\end{minipage} & \begin{minipage}[b]{\linewidth}\raggedright
Valores
\end{minipage} & \begin{minipage}[b]{\linewidth}\raggedright
Significado
\end{minipage} \\
\midrule\noalign{}
\endhead
\bottomrule\noalign{}
\endlastfoot
1º --- natureza & \(\mathrm{K}\) & \textbf{contratual}: EQ + MO. Cobrado
por semana mobilizada, não escala com produção \\
& \(\mathrm{V}\) & \textbf{consumível}: MAT + COMB. Escala com produção
realizada \\
2º --- regime & \(\mathrm{cat}\) & \textbf{catálogo}: nominal,
\(\varepsilon = 1\) \\
& \(\mathrm{op}\) & \textbf{operacional}: com \(\varepsilon_{ij}\)
realizado \\
3º --- markup & \(\beta\) & com BDI aplicado. Ausente = sem BDI \\
\end{longtable}

Exemplos: \(\gamma_i^{\mathrm{K,cat}}\) ·
\(C_{ij}^{\mathrm{V,op},\beta}\) ·
\(\zeta_{ij} = C^{\mathrm{op}}_{ij}/C^{\mathrm{cat}}_{ij}\)

\begin{center}\rule{0.5\linewidth}{0.5pt}\end{center}

\hypertarget{glossuxe1rio-de-suxedmbolos}{%
\subsection{Glossário de Símbolos}\label{glossuxe1rio-de-suxedmbolos}}

\hypertarget{uxedndices-e-conjuntos}{%
\subsubsection{Índices e conjuntos}\label{uxedndices-e-conjuntos}}

\begin{longtable}[]{@{}
  >{\raggedright\arraybackslash}p{(\columnwidth - 2\tabcolsep) * \real{0.5000}}
  >{\raggedright\arraybackslash}p{(\columnwidth - 2\tabcolsep) * \real{0.5000}}@{}}
\toprule\noalign{}
\begin{minipage}[b]{\linewidth}\raggedright
Símbolo
\end{minipage} & \begin{minipage}[b]{\linewidth}\raggedright
Significado
\end{minipage} \\
\midrule\noalign{}
\endhead
\bottomrule\noalign{}
\endlastfoot
\(i\) & índice de atividade \\
\(j\) & índice de semana (calendário) \\
\(r\) & instância de recurso físico \\
\(l\) & categoria de custo indireto \\
\(g\) & tipo físico de bem consumível \\
\(\tau\) & semana relativa ao início da atividade (\(\tau=j-j_i^0\)), ou
nível na decomposição do pool \\
\(a,b\) & semanas de início e fim de um lote de reposição (Parte X) \\
\(\mathcal{A}\) & conjunto de atividades do projeto \\
\(\mathcal{T}\) & conjunto de tipos de recurso \\
\(\mathcal{L}\) & conjunto de categorias de custo indireto \\
\(\mathcal{A}_j\) & conjunto de atividades ativas na semana \(j\) \\
\end{longtable}

\hypertarget{atividade-entrada}{%
\subsubsection{Atividade --- entrada}\label{atividade-entrada}}

\begin{longtable}[]{@{}
  >{\raggedright\arraybackslash}p{(\columnwidth - 2\tabcolsep) * \real{0.5000}}
  >{\raggedright\arraybackslash}p{(\columnwidth - 2\tabcolsep) * \real{0.5000}}@{}}
\toprule\noalign{}
\begin{minipage}[b]{\linewidth}\raggedright
Símbolo
\end{minipage} & \begin{minipage}[b]{\linewidth}\raggedright
Significado
\end{minipage} \\
\midrule\noalign{}
\endhead
\bottomrule\noalign{}
\endlastfoot
\(Q_i\) & quantitativo total da atividade {[}un{]} \\
\(E_i=(\mathcal{G}_i,R_i)\) & equipe: grafo de sinergia + composição de
recursos \\
\(W_i\) & política de mobilização (perfil trapezoidal) \\
\(B_i\) & conjunto de atividades predecessoras \\
\(\theta_i(b)\) & fração de escopo de \(b\) exigida antes do início de
\(A_i\) \\
\((h_i,d_i)\) & horas/dia e dias/semana da jornada da equipe \\
\end{longtable}

\hypertarget{atividade-derivados-de-cronograma}{%
\subsubsection{Atividade --- derivados de
cronograma}\label{atividade-derivados-de-cronograma}}

\begin{longtable}[]{@{}
  >{\raggedright\arraybackslash}p{(\columnwidth - 2\tabcolsep) * \real{0.5000}}
  >{\raggedright\arraybackslash}p{(\columnwidth - 2\tabcolsep) * \real{0.5000}}@{}}
\toprule\noalign{}
\begin{minipage}[b]{\linewidth}\raggedright
Símbolo
\end{minipage} & \begin{minipage}[b]{\linewidth}\raggedright
Significado
\end{minipage} \\
\midrule\noalign{}
\endhead
\bottomrule\noalign{}
\endlastfoot
\(j_i^0\) & semana de início (derivada dos predecessores) \\
\(N_i\) & duração em semanas (derivada por busca monotônica) \\
\(m_{ij}\) & número de frentes mobilizadas de \(A_i\) na semana \(j\) \\
\(\varepsilon_{ij}\) & eficiência semanal da atividade \(i\) na semana
\(j\) \\
\(\varphi_i\) & fração de fechamento da última semana ativa \\
\(M_i\) & total de frente-semanas mobilizadas \\
\(q_{ij}\) & quantitativo produzido por \(A_i\) na semana \(j\) \\
\end{longtable}

\hypertarget{equipe-e-recursos}{%
\subsubsection{Equipe e recursos}\label{equipe-e-recursos}}

\begin{longtable}[]{@{}
  >{\raggedright\arraybackslash}p{(\columnwidth - 2\tabcolsep) * \real{0.5000}}
  >{\raggedright\arraybackslash}p{(\columnwidth - 2\tabcolsep) * \real{0.5000}}@{}}
\toprule\noalign{}
\begin{minipage}[b]{\linewidth}\raggedright
Símbolo
\end{minipage} & \begin{minipage}[b]{\linewidth}\raggedright
Significado
\end{minipage} \\
\midrule\noalign{}
\endhead
\bottomrule\noalign{}
\endlastfoot
\(\mathcal{G}_i=(V_i,A_i^{\text{syn}},\text{cap})\) & grafo de sinergia
da equipe \\
\(D,\ s_i,\ t_i\) & domínio de produção, fonte e sorvedouro virtuais do
grafo \\
\(R_i\) & composição de recursos da equipe (uma frente) \\
\(\text{cat}(r)\) & categoria de BDI do recurso (\(EQ\), \(MO\), \(MAT\)
ou \(COMB\)) \\
\(n_r\) & quantidade do recurso \(r\) por frente \\
\(V_r^{\text{mensal}}\) & custo contratual mensal do recurso (regime
\texttt{mensal}) \\
\(\kappa_r\) & taxa física de consumo do recurso (regime
\texttt{consumo}) \\
\(g(r)\) & tipo físico do bem consumido por \(r\) \\
\(p_{g(r)}\) & preço de aquisição do bem \(g(r)\) \\
\(c_r\) & custo horário de catálogo do recurso \(r\) \\
\(\rho_i\) & produtividade horária da equipe (min-corte de
\(\mathcal{G}_i\)) \\
\(\eta_{s,ij}\) & fator de alocação (grau de utilização do grafo de
sinergia) \\
\(\eta_{p,j},\ \eta_{w,j}\) & fatores globais de desempenho e clima do
calendário \\
\end{longtable}

\hypertarget{mobilizauxe7uxe3o}{%
\subsubsection{Mobilização}\label{mobilizauxe7uxe3o}}

\begin{longtable}[]{@{}
  >{\raggedright\arraybackslash}p{(\columnwidth - 2\tabcolsep) * \real{0.5000}}
  >{\raggedright\arraybackslash}p{(\columnwidth - 2\tabcolsep) * \real{0.5000}}@{}}
\toprule\noalign{}
\begin{minipage}[b]{\linewidth}\raggedright
Símbolo
\end{minipage} & \begin{minipage}[b]{\linewidth}\raggedright
Significado
\end{minipage} \\
\midrule\noalign{}
\endhead
\bottomrule\noalign{}
\endlastfoot
\(m_i^0,\ m_i^1\) & frentes na semana de início e na última semana \\
\(\mu_i^+,\ \mu_i^-\) & taxas de ramp-up e ramp-down
{[}frentes/semana{]} \\
\(m_i^{\max}\) & número máximo de frentes concorrentes (platô) \\
\(m_i(\tau)\) & perfil trapezoidal de frentes em função da semana
relativa \\
\end{longtable}

\hypertarget{paruxe2metros-globais}{%
\subsubsection{Parâmetros globais}\label{paruxe2metros-globais}}

\begin{longtable}[]{@{}
  >{\raggedright\arraybackslash}p{(\columnwidth - 2\tabcolsep) * \real{0.5000}}
  >{\raggedright\arraybackslash}p{(\columnwidth - 2\tabcolsep) * \real{0.5000}}@{}}
\toprule\noalign{}
\begin{minipage}[b]{\linewidth}\raggedright
Símbolo
\end{minipage} & \begin{minipage}[b]{\linewidth}\raggedright
Significado
\end{minipage} \\
\midrule\noalign{}
\endhead
\bottomrule\noalign{}
\endlastfoot
\(s_{\text{mes}}\) & semanas por mês médio (constante de calendário) \\
\(\beta^{EQ},\beta^{MO},\beta^{MAT},\beta^{COMB}\) & taxas de BDI por
categoria de recurso \\
\(\beta^{IND}\) & taxa de BDI para itens indiretos \\
\end{longtable}

\hypertarget{custo-de-catuxe1logo}{%
\subsubsection{Custo de catálogo}\label{custo-de-catuxe1logo}}

\begin{longtable}[]{@{}
  >{\raggedright\arraybackslash}p{(\columnwidth - 2\tabcolsep) * \real{0.5000}}
  >{\raggedright\arraybackslash}p{(\columnwidth - 2\tabcolsep) * \real{0.5000}}@{}}
\toprule\noalign{}
\begin{minipage}[b]{\linewidth}\raggedright
Símbolo
\end{minipage} & \begin{minipage}[b]{\linewidth}\raggedright
Significado
\end{minipage} \\
\midrule\noalign{}
\endhead
\bottomrule\noalign{}
\endlastfoot
\(\gamma_i^{k}\) & custo horário de catálogo da categoria \(k\) na
equipe \(i\) \\
\(\gamma_i^{\mathrm{K,cat}},\ \gamma_i^{\mathrm{V,cat}}\) & blocos
contratual e consumível do custo horário \\
\(\gamma_i^{\mathrm{cat}}\) & custo horário de catálogo total da
equipe \\
\(f_i^{\mathrm V}\) & fração consumível do custo de catálogo \\
\(p_i^{\mathrm{cat}}\) & preço unitário de catálogo {[}R\$/un{]} \\
\(C_i^{\mathrm{cat}}\) & custo de catálogo da atividade {[}R\${]} \\
\end{longtable}

\hypertarget{custos-diretos-operacionais}{%
\subsubsection{Custos diretos
operacionais}\label{custos-diretos-operacionais}}

\begin{longtable}[]{@{}
  >{\raggedright\arraybackslash}p{(\columnwidth - 2\tabcolsep) * \real{0.5000}}
  >{\raggedright\arraybackslash}p{(\columnwidth - 2\tabcolsep) * \real{0.5000}}@{}}
\toprule\noalign{}
\begin{minipage}[b]{\linewidth}\raggedright
Símbolo
\end{minipage} & \begin{minipage}[b]{\linewidth}\raggedright
Significado
\end{minipage} \\
\midrule\noalign{}
\endhead
\bottomrule\noalign{}
\endlastfoot
\(C_{ij}^{\mathrm{K,op}},\ C_{ij}^{\mathrm{V,op}}\) & custo semanal
contratual e consumível (sem BDI) \\
\(C_{ij}^{\mathrm{op}}\) & custo direto semanal total da atividade \\
\(C_i^{\mathrm{K,op}},\ C_i^{\mathrm{V,op}},\ C_i^{\mathrm{op}}\) &
totais por atividade (sem BDI) \\
\(p_i^{\mathrm{op}}\) & custo unitário operacional da atividade
{[}R\$/un{]} \\
\(\zeta_{ij},\ \pi_{ij}\) & razões operacional/catálogo, semanais \\
\(R_i\) & razão operacional/catálogo agregada da atividade \\
\(\Delta_i^{\mathrm{cat}}\) & desvio de catálogo da atividade, em R\$ \\
\end{longtable}

\hypertarget{camada-bdi-por-atividade}{%
\subsubsection{Camada BDI por
atividade}\label{camada-bdi-por-atividade}}

\begin{longtable}[]{@{}
  >{\raggedright\arraybackslash}p{(\columnwidth - 2\tabcolsep) * \real{0.5000}}
  >{\raggedright\arraybackslash}p{(\columnwidth - 2\tabcolsep) * \real{0.5000}}@{}}
\toprule\noalign{}
\begin{minipage}[b]{\linewidth}\raggedright
Símbolo
\end{minipage} & \begin{minipage}[b]{\linewidth}\raggedright
Significado
\end{minipage} \\
\midrule\noalign{}
\endhead
\bottomrule\noalign{}
\endlastfoot
\(\gamma_i^{\mathrm{K,cat},\beta},\ \gamma_i^{\mathrm{V,cat},\beta}\) &
custos horários de catálogo com BDI \\
\(\beta_i^{\mathrm K},\ \beta_i^{\mathrm V}\) & BDI ponderado dos blocos
contratual e consumível \\
\(\delta_{ij},\ \delta_i\) & fração consumível efetiva, semanal e
agregada \\
\(\beta_{ij},\ \beta_i^{\mathrm{eq}}\) & BDI equivalente semanal e por
atividade \\
\(\bar\varepsilon_i\) & eficiência média realizada da atividade (exata,
via fecho de escopo) \\
\(C_i^{D,\beta}\) & custo direto da atividade com BDI \\
\end{longtable}

\hypertarget{pool-de-recursos-contratuais}{%
\subsubsection{Pool de recursos
contratuais}\label{pool-de-recursos-contratuais}}

\begin{longtable}[]{@{}
  >{\raggedright\arraybackslash}p{(\columnwidth - 2\tabcolsep) * \real{0.5000}}
  >{\raggedright\arraybackslash}p{(\columnwidth - 2\tabcolsep) * \real{0.5000}}@{}}
\toprule\noalign{}
\begin{minipage}[b]{\linewidth}\raggedright
Símbolo
\end{minipage} & \begin{minipage}[b]{\linewidth}\raggedright
Significado
\end{minipage} \\
\midrule\noalign{}
\endhead
\bottomrule\noalign{}
\endlastfoot
\(n_j(\tau)\) & demanda semanal de um tipo de recurso \\
\(qt_\tau\) & tamanho de pool necessário para o tipo \(\tau\) \\
\(G_r^*\) & gap de indiferença para perpetuação {[}semanas{]} \\
\(G_r^{\text{viável}}\) & gap máximo absoluto para perpetuação \\
\(T_r^{\text{ativo}},\ T_r^{\text{standby}},\ T_r^{\text{perm}}\) &
semanas ativas, de standby e de permanência total \\
\(f_r^{\text{standby}}\) & fração de standby da instância (diagnóstico
temporal) \\
\(C^{\mathrm{pool}},\ C^{\mathrm{pool},\beta}\) & custo do pool, sem e
com BDI \\
\(C^{MOB,eq}\) & custo de mobilização/desmobilização do pool \\
\end{longtable}

\hypertarget{consumuxedveis-e-estoque}{%
\subsubsection{Consumíveis e estoque}\label{consumuxedveis-e-estoque}}

\begin{longtable}[]{@{}
  >{\raggedright\arraybackslash}p{(\columnwidth - 2\tabcolsep) * \real{0.5000}}
  >{\raggedright\arraybackslash}p{(\columnwidth - 2\tabcolsep) * \real{0.5000}}@{}}
\toprule\noalign{}
\begin{minipage}[b]{\linewidth}\raggedright
Símbolo
\end{minipage} & \begin{minipage}[b]{\linewidth}\raggedright
Significado
\end{minipage} \\
\midrule\noalign{}
\endhead
\bottomrule\noalign{}
\endlastfoot
\(C^{\mathrm{cons}}\) & custo de aquisição de consumíveis (direto) \\
\(\kappa_i^g\) & taxa de consumo agregada do bem \(g\) pela equipe \(i\)
{[}unid\(_g\)/h{]} \\
\(D_{g,ij},\ D_{g,j}\) & demanda física do bem \(g\): por
atividade/semana, e agregada por semana \\
\(U_g\) & unidades físicas do bem \(g\) usadas no projeto inteiro \\
\(S_{g,j}\) & estoque do bem \(g\) ao final da semana \(j\) \\
\(X_{g,j},\ O_{g,j}\) & quantidade pedida e recebida do bem \(g\) na
semana \(j\) \\
\(L_g\) & lead time de entrega do bem \(g\) {[}semanas{]} \\
\(c_g^{\text{hold}},\ c_g^{\text{ped}}\) & custo de posse
{[}R\$/unid\(_g\)/sem{]} e de pedido {[}R\${]} do bem \(g\) \\
\(n_g^{\text{ped}}\) & número de lotes (pedidos) ótimos do bem \(g\) \\
\(f(b),\ \mathrm{Custo}(a,b)\) & custo ótimo acumulado até \(b\), e
custo de um lote \([a,b]\) (Wagner--Whitin) \\
\(C^{\mathrm{inv}}\) & custo de posse de estoque (indireto) \\
\(C^{\mathrm{PED}},\ C^{\mathrm{PED},\beta}\) & custo total de pedidos,
sem e com BDI (indireto) \\
\end{longtable}

\hypertarget{custo-direto-e-indireto-do-projeto}{%
\subsubsection{Custo direto e indireto do
projeto}\label{custo-direto-e-indireto-do-projeto}}

\begin{longtable}[]{@{}ll@{}}
\toprule\noalign{}
Símbolo & Significado \\
\midrule\noalign{}
\endhead
\bottomrule\noalign{}
\endlastfoot
\(C^D,\ C^{D,\beta}\) & custo direto do projeto, sem e com BDI \\
\(n\) & makespan do projeto {[}semanas{]} \\
\(\text{driver}(l)\) & regime de custo indireto (\texttt{perm} ou
\texttt{onetime}) \\
\(C_l^{I}\) & custo indireto por categoria \\
\(C^I,\ C^{I,\beta}\) & custo indireto total, sem e com BDI \\
\end{longtable}

\hypertarget{suxedntese}{%
\subsubsection{Síntese}\label{suxedntese}}

\begin{longtable}[]{@{}
  >{\raggedright\arraybackslash}p{(\columnwidth - 2\tabcolsep) * \real{0.5000}}
  >{\raggedright\arraybackslash}p{(\columnwidth - 2\tabcolsep) * \real{0.5000}}@{}}
\toprule\noalign{}
\begin{minipage}[b]{\linewidth}\raggedright
Símbolo
\end{minipage} & \begin{minipage}[b]{\linewidth}\raggedright
Significado
\end{minipage} \\
\midrule\noalign{}
\endhead
\bottomrule\noalign{}
\endlastfoot
\(C^{\mathrm{cat}}\) & envelope de catálogo do projeto \\
\(C^{\mathrm{tot}}\) & custo operacional total do projeto (sem BDI) \\
\(P\) & preço de venda \\
\(\mathrm{BDI}\) & taxa agregada derivada, \(P/C^{\mathrm{cat}}-1\) \\
\(R\) & razão custo real / catálogo (diagnóstico) \\
\(\beta^{\mathrm{eq}}\) & BDI equivalente sobre o custo total \\
\(\mu\) & margem sobre o preço de venda \\
\(\mathrm{BDI}_{\mathrm{TCU,min}}\) & piso regulatório de BDI (fórmula
TCU) \\
\end{longtable}

\end{document}